% Options for packages loaded elsewhere
\PassOptionsToPackage{unicode}{hyperref}
\PassOptionsToPackage{hyphens}{url}
\documentclass[
  11pt,
]{article}
\usepackage{xcolor}
\usepackage{amsmath,amssymb}
\setcounter{secnumdepth}{-\maxdimen} % remove section numbering
\usepackage{iftex}
\ifPDFTeX
  \usepackage[T1]{fontenc}
  \usepackage[utf8]{inputenc}
  \usepackage{textcomp} % provide euro and other symbols
\else % if luatex or xetex
  \usepackage{unicode-math} % this also loads fontspec
  \defaultfontfeatures{Scale=MatchLowercase}
  \defaultfontfeatures[\rmfamily]{Ligatures=TeX,Scale=1}
\fi
\usepackage{lmodern}
\ifPDFTeX\else
  % xetex/luatex font selection
\fi
% Use upquote if available, for straight quotes in verbatim environments
\IfFileExists{upquote.sty}{\usepackage{upquote}}{}
\IfFileExists{microtype.sty}{% use microtype if available
  \usepackage[]{microtype}
  \UseMicrotypeSet[protrusion]{basicmath} % disable protrusion for tt fonts
}{}
\makeatletter
\@ifundefined{KOMAClassName}{% if non-KOMA class
  \IfFileExists{parskip.sty}{%
    \usepackage{parskip}
  }{% else
    \setlength{\parindent}{0pt}
    \setlength{\parskip}{6pt plus 2pt minus 1pt}}
}{% if KOMA class
  \KOMAoptions{parskip=half}}
\makeatother
\usepackage{color}
\usepackage{fancyvrb}
\newcommand{\VerbBar}{|}
\newcommand{\VERB}{\Verb[commandchars=\\\{\}]}
\DefineVerbatimEnvironment{Highlighting}{Verbatim}{commandchars=\\\{\}}
% Add ',fontsize=\small' for more characters per line
\newenvironment{Shaded}{}{}
\newcommand{\AlertTok}[1]{\textcolor[rgb]{1.00,0.00,0.00}{\textbf{#1}}}
\newcommand{\AnnotationTok}[1]{\textcolor[rgb]{0.38,0.63,0.69}{\textbf{\textit{#1}}}}
\newcommand{\AttributeTok}[1]{\textcolor[rgb]{0.49,0.56,0.16}{#1}}
\newcommand{\BaseNTok}[1]{\textcolor[rgb]{0.25,0.63,0.44}{#1}}
\newcommand{\BuiltInTok}[1]{\textcolor[rgb]{0.00,0.50,0.00}{#1}}
\newcommand{\CharTok}[1]{\textcolor[rgb]{0.25,0.44,0.63}{#1}}
\newcommand{\CommentTok}[1]{\textcolor[rgb]{0.38,0.63,0.69}{\textit{#1}}}
\newcommand{\CommentVarTok}[1]{\textcolor[rgb]{0.38,0.63,0.69}{\textbf{\textit{#1}}}}
\newcommand{\ConstantTok}[1]{\textcolor[rgb]{0.53,0.00,0.00}{#1}}
\newcommand{\ControlFlowTok}[1]{\textcolor[rgb]{0.00,0.44,0.13}{\textbf{#1}}}
\newcommand{\DataTypeTok}[1]{\textcolor[rgb]{0.56,0.13,0.00}{#1}}
\newcommand{\DecValTok}[1]{\textcolor[rgb]{0.25,0.63,0.44}{#1}}
\newcommand{\DocumentationTok}[1]{\textcolor[rgb]{0.73,0.13,0.13}{\textit{#1}}}
\newcommand{\ErrorTok}[1]{\textcolor[rgb]{1.00,0.00,0.00}{\textbf{#1}}}
\newcommand{\ExtensionTok}[1]{#1}
\newcommand{\FloatTok}[1]{\textcolor[rgb]{0.25,0.63,0.44}{#1}}
\newcommand{\FunctionTok}[1]{\textcolor[rgb]{0.02,0.16,0.49}{#1}}
\newcommand{\ImportTok}[1]{\textcolor[rgb]{0.00,0.50,0.00}{\textbf{#1}}}
\newcommand{\InformationTok}[1]{\textcolor[rgb]{0.38,0.63,0.69}{\textbf{\textit{#1}}}}
\newcommand{\KeywordTok}[1]{\textcolor[rgb]{0.00,0.44,0.13}{\textbf{#1}}}
\newcommand{\NormalTok}[1]{#1}
\newcommand{\OperatorTok}[1]{\textcolor[rgb]{0.40,0.40,0.40}{#1}}
\newcommand{\OtherTok}[1]{\textcolor[rgb]{0.00,0.44,0.13}{#1}}
\newcommand{\PreprocessorTok}[1]{\textcolor[rgb]{0.74,0.48,0.00}{#1}}
\newcommand{\RegionMarkerTok}[1]{#1}
\newcommand{\SpecialCharTok}[1]{\textcolor[rgb]{0.25,0.44,0.63}{#1}}
\newcommand{\SpecialStringTok}[1]{\textcolor[rgb]{0.73,0.40,0.53}{#1}}
\newcommand{\StringTok}[1]{\textcolor[rgb]{0.25,0.44,0.63}{#1}}
\newcommand{\VariableTok}[1]{\textcolor[rgb]{0.10,0.09,0.49}{#1}}
\newcommand{\VerbatimStringTok}[1]{\textcolor[rgb]{0.25,0.44,0.63}{#1}}
\newcommand{\WarningTok}[1]{\textcolor[rgb]{0.38,0.63,0.69}{\textbf{\textit{#1}}}}
\usepackage{longtable,booktabs,array}
\newcounter{none} % for unnumbered tables
\usepackage{calc} % for calculating minipage widths
% Correct order of tables after \paragraph or \subparagraph
\usepackage{etoolbox}
\makeatletter
\patchcmd\longtable{\par}{\if@noskipsec\mbox{}\fi\par}{}{}
\makeatother
% Allow footnotes in longtable head/foot
\IfFileExists{footnotehyper.sty}{\usepackage{footnotehyper}}{\usepackage{footnote}}
\makesavenoteenv{longtable}
\setlength{\emergencystretch}{3em} % prevent overfull lines
\providecommand{\tightlist}{%
  \setlength{\itemsep}{0pt}\setlength{\parskip}{0pt}}
% Shared preamble for the LAMMPS-CO manuscript, supplementary, and cover letter.
% Compiled with xelatex. Latin Modern is used rather than a system font so the
% build is reproducible on any TeX installation; the characters it lacks
% (Greek, super/subscripts, arrows, a few symbols) are mapped explicitly below.
\usepackage{fontspec}
\setmainfont{Latin Modern Roman}
\setsansfont{Latin Modern Sans}
\setmonofont{Latin Modern Mono}[Scale=0.85]

\usepackage[a4paper,margin=1in]{geometry}
\usepackage{longtable,booktabs,array}
\usepackage{graphicx}
\usepackage{amsmath,amssymb}
\usepackage[table]{xcolor}
\usepackage[hidelinks,breaklinks=true]{hyperref}
\usepackage{microtype}
\usepackage{newunicodechar}

% Characters that appear in the text but not in Latin Modern's text font.
\newunicodechar{λ}{\ensuremath{\lambda}}
\newunicodechar{Δ}{\ensuremath{\Delta}}
\newunicodechar{ε}{\ensuremath{\varepsilon}}
\newunicodechar{σ}{\ensuremath{\sigma}}
\newunicodechar{ρ}{\ensuremath{\rho}}
\newunicodechar{τ}{\ensuremath{\tau}}
\newunicodechar{φ}{\ensuremath{\varphi}}
\newunicodechar{δ}{\ensuremath{\delta}}
\newunicodechar{α}{\ensuremath{\alpha}}
\newunicodechar{θ}{\ensuremath{\theta}}
\newunicodechar{→}{\ensuremath{\rightarrow}}
\newunicodechar{≈}{\ensuremath{\approx}}
\newunicodechar{√}{\ensuremath{\sqrt{}}}
\newunicodechar{∫}{\ensuremath{\int}}
\newunicodechar{⟨}{\ensuremath{\langle}}
\newunicodechar{⟩}{\ensuremath{\rangle}}
\newunicodechar{■}{\ensuremath{\blacksquare}}
\newunicodechar{⁺}{\ensuremath{^{+}}}
\newunicodechar{⁻}{\ensuremath{^{-}}}
\newunicodechar{⁰}{\ensuremath{^{0}}}
\newunicodechar{⁴}{\ensuremath{^{4}}}
\newunicodechar{⁵}{\ensuremath{^{5}}}
\newunicodechar{⁶}{\ensuremath{^{6}}}
\newunicodechar{⁷}{\ensuremath{^{7}}}
\newunicodechar{₀}{\ensuremath{_{0}}}
\newunicodechar{₁}{\ensuremath{_{1}}}
\newunicodechar{₂}{\ensuremath{_{2}}}
\newunicodechar{₃}{\ensuremath{_{3}}}
\newunicodechar{₄}{\ensuremath{_{4}}}
\newunicodechar{₆}{\ensuremath{_{6}}}

% pandoc emits this when a document mixes tables and lists
\providecommand{\tightlist}{%
  \setlength{\itemsep}{0pt}\setlength{\parskip}{0pt}}

% Long tables of file names need to break across pages
\setlength{\LTleft}{0pt}
\setlength{\LTright}{0pt}
\renewcommand{\arraystretch}{1.15}

\setlength{\parskip}{0.6em}
\setlength{\parindent}{0pt}
\sloppy
\usepackage{bookmark}
\IfFileExists{xurl.sty}{\usepackage{xurl}}{} % add URL line breaks if available
\urlstyle{same}
% fallback for those not using the hyperref driver hyperxmp:
\makeatletter
\@ifundefined{xmpquote}{\newcommand{\xmpquote}[1]{#1}}{}
\makeatother
\hypersetup{
  hidelinks,
  pdfcreator={LaTeX via pandoc}}

\author{}
\date{}

\begin{document}

\section{LAMMPS-CO: extended OpenMP coverage and compiler-friendly force
kernels for 261 LAMMPS pair
styles}\label{lammps-co-extended-openmp-coverage-and-compiler-friendly-force-kernels-for-261-lammps-pair-styles}

\textbf{Chun-Yeol You¹}

¹ Department of Physics and Chemistry, DGIST (Daegu Gyeongbuk Institute
of Science and Technology), Daegu, Republic of Korea

\emph{Correspondence: cyyou@dgist.ac.kr}

\begin{center}\rule{0.5\linewidth}{0.5pt}\end{center}

\subsection{Abstract}\label{abstract}

LAMMPS ships 307 pair interaction styles, but only 157 of them have an
OpenMP-accelerated variant, and the serial force kernels of the
remainder are written in a form that inhibits compiler
auto-vectorization: positions and forces are reached through unqualified
\texttt{double\ **} pointers, and forces are accumulated by repeated
read--modify--write of the per-atom array inside the neighbor loop.
LAMMPS-CO supplies 36 new OpenMP pair-style variants, raising coverage
to 193 of 307 styles (62.9\%), and rewrites the force kernels of 225
further pair files across 38 packages to use
\texttt{\_\_restrict\_\_}-qualified typed pointers and local
accumulators --- the same patterns the OPENMP and OPT packages already
rely on. Redundant \texttt{pow()} evaluations are eliminated from the
\texttt{nm/cut}, \texttt{lj/pirani}, and \texttt{mie/cut} families. On a
32,000-atom Lennard-Jones benchmark the new OpenMP variants reach 4.31×
at four threads and 4.93× at eight; the kernel rewrite yields 4--8\% on
pair time, the compiler-flag configuration a further 6.3\%, and the
\texttt{pow()} elimination 19.9--40.8\% within the affected styles.
Correctness is established by running the unmodified and modified builds
of the same LAMMPS commit through the upstream
\texttt{unittest/force-styles} regression suite --- 325 pair-style
tests, identical outcomes --- supplemented by dedicated decks,
distributed with the code, for the eleven styles that suite does not
reach. The newly parallelized styles span free-energy-perturbation
soft-core potentials, dielectric continuum models, and
dispersion-corrected electrostatics; as a worked application we compute
Li⁺ desolvation free energies in ethylene carbonate, propylene
carbonate, and 1,2-dimethoxyethane by OPLS-AA thermodynamic integration,
a 63-window campaign that the OpenMP variants bring within reach of a
single workstation. The code was produced with the assistance of a large
language model coding agent, and the generation protocol, its measured
error rate, and the verification that caught those errors are reported
alongside the software.

\textbf{Keywords}: molecular dynamics, LAMMPS, OpenMP parallelization,
force kernel optimization, auto-vectorization, free energy perturbation,
battery electrolyte, LLM-assisted code generation

\begin{center}\rule{0.5\linewidth}{0.5pt}\end{center}

\subsection{Program summary}\label{program-summary}

\emph{Program Title:} LAMMPS-CO

\emph{CPC Library link to program files:} (to be added by Technical
Editor)

\emph{Developer's repository link:}
\url{https://github.com/mirryou-maker/lammps-CO}

\emph{Licensing provisions:} GPLv2

\emph{Programming language:} C++ (C++11), OpenMP

\emph{Supplementary material:} file-level inventory of all modified pair
styles (Supplementary Table S6); prompt templates used for code
generation (Supplementary Note S1)

\emph{Nature of problem:} LAMMPS parallelizes across MPI ranks by
spatial decomposition, but within a rank a pair style is threaded only
if an OpenMP variant of that style exists. Fewer than half of the 307
pair styles have one, so on a many-core node the untreated styles leave
the cores idle. Independently, the serial force kernels are written in a
form the compiler cannot vectorize: aliasing between the position and
force arrays cannot be ruled out, and the per-atom force array is
updated inside the innermost loop, creating a loop-carried dependence.

\emph{Solution method:} For each style lacking one, an OpenMP variant is
generated following the structure of the existing OPENMP package: the
neighbor-list loop is partitioned across threads, per-thread force and
virial accumulation is routed through \texttt{ThrData}, and the
Newton-pair and energy/virial tallying paths are preserved. For the
serial kernels, positions and forces are reinterpreted through a
\texttt{\_\_restrict\_\_}-qualified three-component struct pointer, and
the innermost force accumulation is moved into scalar locals that are
written back once per outer iteration. Neither transformation changes
the order of floating-point operations, so results are unchanged; this
is verified rather than assumed, against the reference values in
LAMMPS's own force-style regression suite.

\emph{Additional comments including restrictions and unusual features:}
The package is applied to an existing LAMMPS source tree by
\texttt{scripts/apply\_patch.sh}, which walks this repository's
\texttt{src/} and replaces matching files in the target tree; no CMake
changes are needed, as LAMMPS discovers \texttt{\_omp.cpp} files
automatically. Developed against LAMMPS \texttt{91d4111a9} (develop, 5
June 2026). Three files were reverted after the rewrite proved not to
preserve operation order for them (\texttt{pair\_buck\_coul\_cut},
\texttt{pair\_kolmogorov\_crespi\_z}, \texttt{pair\_polymorphic}). The
verification harness in \texttt{tools/verify/} requires two builds of
LAMMPS --- one unmodified, one with the package applied --- from the
same commit.

\begin{center}\rule{0.5\linewidth}{0.5pt}\end{center}

\subsection{Introduction}\label{introduction}

Molecular dynamics simulation has transformed materials research over
the past three decades, providing atom-level insight that complements
and guides experiment across scales from quantum defect cores to
mesoscale grain boundaries {[}1,2{]}. The scope of MD application in
materials science has expanded dramatically: from classical force-field
studies of fracture mechanics {[}3{]} and ionic conductivity {[}4,5{]}
to machine-learning interatomic potentials (MLIPs) that approach
density-functional theory accuracy at a fraction of the computational
cost {[}6,7,8{]}. As simulation systems grow toward billions of atoms
and microsecond timescales on current leadership-class supercomputers,
the efficiency of the underlying MD engine becomes a scientific
bottleneck as much as a computational one {[}9{]}.

LAMMPS {[}1,2,10,11{]}, developed at Sandia National Laboratories and
released under GPLv2, is the de facto standard MD engine in materials
science with over 10,000 citations of its principal reference paper
{[}2{]}. Its plugin-style architecture---more than 300 pair interaction
styles, 100 packages, and a rich scripting interface---enables
simulation of systems ranging from atomic-scale metallic alloys {[}12{]}
to coarse-grained polymer networks {[}13{]} and reactive chemistry
{[}14,18{]}. Critically for materials science, LAMMPS has been the
computational engine behind a series of landmark discoveries:
dislocation-mediated high-temperature strength in refractory BCC alloys
{[}15{]}, dynamic nanodomains in hybrid perovskites {[}16{]}, and
superionic paddle-wheel mechanisms in sodium halide electrolytes
{[}17{]}.

Despite this maturity, LAMMPS retains well-documented optimization gaps
in its standard (non-accelerated) code path. The OpenMP acceleration
package (OPENMP), which delivers 2--8× speedup through thread-level
parallelism via a thread-safe force-accumulation design {[}1{]}, covers
only 157 of 307 pair\_style implementations (51\%), leaving
approximately 150 pair styles---widely used in FEP alchemical
calculations, dielectric media simulations, and dispersion-corrected
electrostatics---without thread-parallel variants. In parallel, the
serial hotloop in many standard pair files contains unoptimized pointer
patterns that prevent compiler auto-vectorization, a limitation that the
OPT package {[}1{]} addresses for only 15 pair styles through
template-based branch elimination.

Bridging this gap through manual coding would require hundreds of
developer-hours and is prone to thread-safety errors that are difficult
to detect without exhaustive numerical validation. The recent emergence
of LLM-based coding agents---systems capable of reading, understanding,
and transforming large C++ codebases---offers a new pathway for
systematic, pattern-based code generation at scale {[}19,24,31{]}.
Claude Code (Anthropic), a terminal-native LLM coding agent, combines
static analysis of the codebase with iterative code generation and can
be guided through complex multi-file transformations via
natural-language prompts.

In this work, we demonstrate the first systematic use of an LLM coding
agent to extend OpenMP parallelization coverage and apply
compiler-friendly micro-optimizations across the LAMMPS codebase. Our
contributions are: (i) a six-step, risk-stratified optimization protocol
applicable to any large scientific software project; (ii) 36 new OMP
pair\_style variants, verified bit-identically against their serial
counterparts; (iii) restrict-annotated local-accumulator
(\texttt{fxtmp}) patterns applied to 225 standard pair files; (iv) a
reusable automated benchmarking harness for regression testing; and (v)
quantified speedup data covering representative simulation workloads
from battery electrolyte modeling to free-energy perturbation.

\begin{center}\rule{0.5\linewidth}{0.5pt}\end{center}

\subsection{Results}\label{results}

\subsubsection{Overview of the Optimization
Protocol}\label{overview-of-the-optimization-protocol}

The optimization workflow (Figure 1) was structured into six steps
ordered by increasing code-change risk and decreasing suitability for
LLM-pattern replication:

\begin{enumerate}
\def\labelenumi{\arabic{enumi}.}
\tightlist
\item
  \textbf{Step B-1} --- Automated benchmarking infrastructure
\item
  \textbf{Step B-2} --- Acceleration package coverage diagnostics
\item
  \textbf{Step A-2} --- Compiler flag tuning (no source changes)
\item
  \textbf{Step A-1} --- OpenMP pair\_style variant generation (new
  files)
\item
  \textbf{Step A-3} --- restrict/fxtmp backport to standard serial
  hotloops (file modification)
\item
  \textbf{Step A-4} --- Transcendental function deduplication (targeted
  pow() elimination)
\end{enumerate}

Each step was preceded by a structured Claude Code prompt specifying the
target files, the exact transformation pattern from an existing
reference implementation, and the numerical verification criterion. Step
A-1 and A-3 prompts followed a canonical format: (1) identify the gap,
(2) state the reference implementation (e.g.,
\texttt{src/OPENMP/pair\_lj\_cut\_omp.cpp}), (3) describe the
transformation rules, (4) specify the verification test. Typical prompt
length was 300--600 tokens; the codebase context (relevant files) was
loaded via explicit file-read instructions. The full prompt templates
used in this work are provided in Supplementary Methods.

\begin{center}\rule{0.5\linewidth}{0.5pt}\end{center}

\textbf{{[}FIGURE 1{]}}\\
\emph{Figure 1: AI-assisted LAMMPS optimization workflow. (a) Six-step
optimization pipeline organized by risk and LLM suitability. Risk
increases from Step B-1 (no code changes) to Step A-3 (in-place
modification of 228 files). (b) Claude Code interaction pattern: the LLM
reads reference implementations, extracts transformation patterns,
generates new code, and triggers automated numerical verification. (c)
The three-level verification pyramid: build-time (compilation),
single-step (bit-identical thermo output), and multi-step (NVE energy
conservation).}

\begin{center}\rule{0.5\linewidth}{0.5pt}\end{center}

\subsubsection{Benchmarking Infrastructure (Step
B-1)}\label{benchmarking-infrastructure-step-b-1}

To enable objective, reproducible performance measurement throughout the
project, we first constructed an automated benchmarking harness
(\texttt{tools/bench/Run-Benchmark.ps1}). The harness executes LAMMPS
with a specified input file N times, parses the MPI timing breakdown
from stdout, and stores results in structured JSON and a cumulative CSV
database. Two canonical benchmark inputs were defined: a small-system
case (4,000 atoms, FCC Ar, \texttt{lj/cut}, NVE, 250 steps) for fast
iteration, and a large-system case (32,000 atoms, same conditions, 500
steps) providing \textasciitilde5\% coefficient of variation on loop
time and serving as the primary reference for speedup claims (Table 1).

\begin{center}\rule{0.5\linewidth}{0.5pt}\end{center}

\textbf{{[}TABLE 1{]}}\\
\emph{Table 1: Baseline performance measurements. Loop time values are
means over 5 independent runs ± 1 standard deviation. Pair computation
accounts for 80.5--80.8\% of total loop time in both cases.}

{\def\LTcaptype{none} % do not increment counter
\begin{longtable}[]{@{}llllll@{}}
\toprule\noalign{}
System & Atoms & Steps & Loop time (s) & Std (s) & Pair fraction \\
\midrule\noalign{}
\endhead
\bottomrule\noalign{}
\endlastfoot
Small & 4,000 & 250 & 0.289 & 0.041 & 80.5\% \\
Large & 32,000 & 500 & 4.333 & 0.112 & 80.8\% \\
\end{longtable}
}

\begin{center}\rule{0.5\linewidth}{0.5pt}\end{center}

The harness also includes per-pair-style correctness test inputs
(\texttt{in.*\_test}), which run 50 NVE steps with high-precision thermo
output (\texttt{thermo\_modify\ format\ float\ \%20.15g}) and compare
serial vs.~OMP results at the last significant digit.

\subsubsection{Acceleration Package Coverage and Diagnostics (Step
B-2)}\label{acceleration-package-coverage-and-diagnostics-step-b-2}

Analysis of the LAMMPS source tree revealed significant gaps in OpenMP
acceleration coverage (Figure 2a). Of 307 unique \texttt{pair\_style}
implementations, only 157 (51.1\%) had OMP variants in
\texttt{src/OPENMP/}. Coverage was higher for bonded interactions
(57--81\%) but remained incomplete. The INTEL and OPT acceleration
packages covered even fewer styles (19 and 15 respectively),
highlighting that the OMP package represents the most impactful
near-term target.

A PowerShell diagnostic script
(\texttt{tools/accel/Check-AccelCoverage.ps1}) was written to parse user
input scripts and flag styles for which OMP variants exist but are not
activated, providing actionable recommendations
(\texttt{-sf\ omp\ -pk\ omp\ N} flags) without any code modification.
This tool addresses the deployment gap: many LAMMPS users run on
multi-core workstations without activating the OPENMP package,
forfeiting 3--5× available speedup.

\begin{center}\rule{0.5\linewidth}{0.5pt}\end{center}

\textbf{{[}FIGURE 2{]}}\\
\emph{Figure 2: OpenMP acceleration coverage in LAMMPS before and after
optimization. (a) Coverage fractions for each interaction category
before optimization (grey) and after addition of 36 new pair\_style OMP
variants (colored). Error bars indicate ±1 for single-style uncertainty
in manual counting. (b) Breakdown of the 36 new OMP pair variants by
source package. (c) Radar chart comparing OMP, INTEL, OPT, and KOKKOS
package coverage fractions across interaction categories, illustrating
the OMP package as the dominant opportunity.}

\begin{center}\rule{0.5\linewidth}{0.5pt}\end{center}

\subsubsection{Compiler Flag Optimization (Step
A-2)}\label{compiler-flag-optimization-step-a-2}

The default LAMMPS build on Windows uses MSVC with
\texttt{CMAKE\_BUILD\_TYPE=Release} and no additional CPU-specific
flags. We evaluated three flag combinations against this baseline using
the large-system benchmark (Table 2, Figure 3). Enabling
\texttt{/arch:AVX2} produced a consistent 2.9\% improvement in loop time
on the large system (4.209 s vs.~4.333 s baseline), attributable to
256-bit SIMD vectorization of the inner neighbor-list loop. Adding
\texttt{/fp:fast} in isolation yielded a similar gain (+2.8\%), but the
combination \texttt{/arch:AVX2\ /fp:fast} achieved a synergistic 6.3\%
improvement (4.058 s), consistent with the compiler exploiting fused
multiply-add (FMA) instructions once IEEE rounding constraints are
relaxed. LTO (\texttt{/GL+LTCG}) showed marginal gain (+2.1\%) with high
run-to-run variance and substantially increased link time; we do not
recommend it for routine LAMMPS builds.

\begin{center}\rule{0.5\linewidth}{0.5pt}\end{center}

\textbf{{[}FIGURE 3{]}}\\
\emph{Figure 3: Compiler flag optimization on the 32,000-atom benchmark.
(a) Loop time (mean ± 1 std, N=5 runs) for baseline and four flag
configurations. Horizontal dashed line indicates baseline. (b)
Percentage improvement in loop time relative to baseline. The
AVX2+/fp:fast combination delivers 6.3\% improvement while LTO shows
marginal and variable gain. (c) Pair time fraction vs.~configuration,
confirming that non-pair sections (neighbor, comm) are not adversely
affected.}

\begin{center}\rule{0.5\linewidth}{0.5pt}\end{center}

\textbf{{[}TABLE 2{]}}\\
\emph{Table 2: Build flag optimization results on 32,000-atom FCC
Lennard-Jones benchmark (500 steps, serial execution, N=5 runs). Speedup
is relative to the no-flag baseline.}

{\def\LTcaptype{none} % do not increment counter
\begin{longtable}[]{@{}llll@{}}
\toprule\noalign{}
Configuration & Loop time (s) & Std (s) & Speedup vs.~baseline \\
\midrule\noalign{}
\endhead
\bottomrule\noalign{}
\endlastfoot
Baseline (no extra flags) & 4.333 & 0.112 & 1.000× \\
\texttt{/arch:AVX2} & 4.209 & 0.169 & 1.029× \\
\texttt{/fp:fast} & 4.214 & 0.043 & 1.028× \\
LTO (\texttt{/GL+LTCG}) & 4.245 & 0.199 & 1.021× \\
\texttt{/arch:AVX2} + \texttt{/fp:fast} & 4.058 & 0.124 & 1.068× \\
\end{longtable}
}

\begin{quote}
\textbf{Note on \texttt{/fp:fast}}: This flag disables strict IEEE 754
compliance (associativity relaxation, unsafe math optimizations). For
materials simulation, where energy conservation is the primary
correctness criterion, \texttt{/fp:fast} is acceptable as an opt-in
performance preset. We observed no change in thermodynamic trajectories
at the 5-decimal-place level in 50-step NVE validation runs.
\end{quote}

\subsubsection{OpenMP Pair\_Style Coverage Extension (Step
A-1)}\label{openmp-pair_style-coverage-extension-step-a-1}

\paragraph{Code Generation
Methodology}\label{code-generation-methodology}

The LAMMPS OPENMP package follows a well-documented template pattern
described in \texttt{doc/src/Developer\_write\_openmp.rst}: a new
\texttt{PairXxxOMP} class inherits from both \texttt{PairXxx} (the
standard implementation) and \texttt{ThrOMP} (the thread-safety
infrastructure), overrides \texttt{compute()} with an OpenMP parallel
region that dispatches to a templated
\texttt{eval\textless{}EVFLAG,EFLAG\textgreater{}()} method, and uses
\texttt{thr-\textgreater{}get\_f()} per-thread force buffers with
\texttt{ev\_tally\_thr()} for thread-safe energy/virial accumulation
(Figure 4a).

Claude Code was prompted with: (1) the target style's \texttt{.h} and
\texttt{.cpp} files, (2) the closest existing OMP variant as a
reference, (3) explicit rules for the three structurally distinct cases
(half-list with Newton, half-list without Newton, and full-list for
DIELECTRIC styles with electric field accumulation), and (4) the
bit-identical verification test. For each generated file, Claude Code
also produced the corresponding CMakeLists entry and package
registration line.

\paragraph{New OMP Variants}\label{new-omp-variants}

Thirty-five new OMP pair variants were generated and verified across
four package groups (Table 3):

\begin{itemize}
\tightlist
\item
  \textbf{EXTRA-PAIR (18 styles)}: FEP soft-core potentials
  (\texttt{lj/cut/soft}, \texttt{lj/cut/coul/long/soft},
  \texttt{morse/soft}, \texttt{coul/cut/soft/gapsys}, etc.) and
  specialty potentials (\texttt{nm/cut/split}, \texttt{born/coul/dsf},
  \texttt{born/coul/wolf/cs}, \texttt{dpd/fdt}, \texttt{bpm/spring}).
\item
  \textbf{DIELECTRIC (7 styles)}: Continuum dielectric boundary models
  with electric field accumulation (\texttt{lj/cut/coul/cut/dielectric},
  \texttt{lj/cut/coul/long/dielectric}, \texttt{coul/long/dielectric},
  etc.), requiring specialized full-list treatment with per-atom ε
  arrays.
\item
  \textbf{KSPACE/long-range (3 styles)}: Dispersion long-range
  potentials with Ewald summation (\texttt{lj/long/coul/long},
  \texttt{buck/long/coul/long}, \texttt{lj/long/tip4p/long}), requiring
  co-dispatch of real-space and k-space contributions.
\item
  \textbf{Other packages (7 styles)}: Styles from DPD-MESO
  (\texttt{dpd/fdt/energy}, \texttt{mdpd}), INTERLAYER
  (\texttt{kolmogorov/crespi/z}, \texttt{lebedeva/z}), DRUDE
  (\texttt{thole}), RHEO (\texttt{rheo/solid}), and DIPOLE
  (\texttt{lj/cut/dipole/long}) packages, each following the standard
  half-list OMP template with no additional structural modifications.
\end{itemize}

\begin{center}\rule{0.5\linewidth}{0.5pt}\end{center}

\textbf{{[}TABLE 3{]}}\\
\emph{Table 3: Newly added OpenMP pair\_style variants. The OMP variant
column lists the new styles; all are bit-identical to their serial
counterparts on the 50-step NVE validation test.}

{\def\LTcaptype{none} % do not increment counter
\begin{longtable}[]{@{}
  >{\raggedright\arraybackslash}p{(\linewidth - 6\tabcolsep) * \real{0.1552}}
  >{\raggedright\arraybackslash}p{(\linewidth - 6\tabcolsep) * \real{0.2586}}
  >{\raggedright\arraybackslash}p{(\linewidth - 6\tabcolsep) * \real{0.2931}}
  >{\raggedright\arraybackslash}p{(\linewidth - 6\tabcolsep) * \real{0.2931}}@{}}
\toprule\noalign{}
\begin{minipage}[b]{\linewidth}\raggedright
Package
\end{minipage} & \begin{minipage}[b]{\linewidth}\raggedright
Standard style
\end{minipage} & \begin{minipage}[b]{\linewidth}\raggedright
New OMP variant
\end{minipage} & \begin{minipage}[b]{\linewidth}\raggedright
Physical context
\end{minipage} \\
\midrule\noalign{}
\endhead
\bottomrule\noalign{}
\endlastfoot
EXTRA-PAIR & \texttt{lj/cut/soft} & \texttt{lj/cut/soft/omp} & FEP
alchemical transformation \\
EXTRA-PAIR & \texttt{lj/cut/coul/long/soft} &
\texttt{lj/cut/coul/long/soft/omp} & FEP with long-range Coulomb \\
EXTRA-PAIR & \texttt{lj/class2/soft} & \texttt{lj/class2/soft/omp} &
Class2 FF with soft core \\
EXTRA-PAIR & \texttt{morse/soft} & \texttt{morse/soft/omp} & Soft-core
Morse potential \\
EXTRA-PAIR & \texttt{coul/cut/soft/gapsys} &
\texttt{coul/cut/soft/gapsys/omp} & Gapsys soft-core electrostatics \\
EXTRA-PAIR & \texttt{nm/cut/split} & \texttt{nm/cut/split/omp} & n-m
Lennard-Jones, split form \\
EXTRA-PAIR & \texttt{born/coul/dsf} & \texttt{born/coul/dsf/omp} & Born
ionic + damped shifted Coulomb \\
EXTRA-PAIR & \texttt{born/coul/wolf/cs} & \texttt{born/coul/wolf/cs/omp}
& Born + Wolf with core-shell \\
EXTRA-PAIR & \texttt{bpm/spring} & \texttt{bpm/spring/omp} &
Bond-particle mechanics \\
EXTRA-PAIR & \texttt{dpd/fdt} & \texttt{dpd/fdt/omp} &
Fluctuation-dissipation DPD \\
EXTRA-PAIR & \texttt{dpd/fdt/energy} & \texttt{dpd/fdt/energy/omp} &
Energy-conserving DPD \\
DIELECTRIC & \texttt{lj/cut/coul/cut/dielectric} &
\texttt{lj/cut/coul/cut/dielectric/omp} & Implicit solvent LJ+Coulomb \\
DIELECTRIC & \texttt{lj/cut/coul/long/dielectric} &
\texttt{lj/cut/coul/long/dielectric/omp} & Implicit solvent + PPPM \\
DIELECTRIC & \texttt{coul/long/dielectric} &
\texttt{coul/long/dielectric/omp} & Pure Coulomb dielectric \\
DIELECTRIC & \texttt{lj/long/coul/long/dielectric} &
\texttt{lj/long/coul/long/dielectric/omp} & Dispersion+Coulomb
dielectric \\
KSPACE & \texttt{lj/long/coul/long} & \texttt{lj/long/coul/long/omp} &
LJ dispersion + PPPM \\
KSPACE & \texttt{buck/long/coul/long} & \texttt{buck/long/coul/long/omp}
& Buckingham + PPPM \\
KSPACE & \texttt{lj/long/tip4p/long} & \texttt{lj/long/tip4p/long/omp} &
TIP4P water + dispersion \\
\end{longtable}
}

\emph{(17 additional styles listed in Supplementary Table S1)}

\paragraph{Benchmarking of New OMP
Styles}\label{benchmarking-of-new-omp-styles}

Three representative newly parallelized styles were benchmarked (Figure
4b, Table 4). The \texttt{born/coul/dsf} style---relevant to oxide and
ionic crystal simulation---achieved a 2.92× speedup at 4 OMP threads and
3.96× at 8 threads; \texttt{lj/class2/soft}, used in Class-2 FEP
calculations for biomolecular free energies, achieved 3.14× (4t) and
3.75× (8t); \texttt{nm/cut/split} achieved 2.63× (4t) and 2.81× (8t).
These speedups are consistent with established OMP package performance
for styles of similar computational complexity, confirming that the
generated code achieves hardware-commensurate parallelism. The lower
speedup of \texttt{nm/cut/split} relative to the other two styles
reflects its higher Amdahl serial fraction (f = 17\%) driven by
neighbor-list building overhead (\textasciitilde10\% of loop time),
which does not benefit from OMP pair parallelism.

On the primary 32,000-atom \texttt{lj/cut} benchmark, the OMP package
(\texttt{-sf\ omp\ -pk\ omp\ N}) delivers the performance profile shown
in Figure 4c and Table 5. At 4 threads, loop time decreases from 4.333 s
to 1.004 s (4.31× speedup); at 8 threads, 0.878 s (4.93×). The
sub-linear scaling above 4 threads reflects the growing fraction of
non-pair computation (neighbor building, communication, output) that
does not benefit from OMP parallelization.

\begin{center}\rule{0.5\linewidth}{0.5pt}\end{center}

\textbf{{[}FIGURE 4{]}}\\
\emph{Figure 4: OpenMP pair\_style optimization. (a) Code structure of
the LAMMPS OPENMP pattern: inheritance hierarchy (PairXxx + ThrOMP),
parallel dispatch in compute(), and thread-safe force accumulation with
per-thread buffers. (b) Speedup at 4 OMP threads for three
representative newly parallelized styles (864 atoms, 500 NVE steps, N=3
runs). Error bars are ±1σ. (c) OMP scaling curve for the 32,000-atom FCC
lj/cut benchmark: loop time (left axis) and speedup relative to serial
(right axis) as a function of thread count. Ideal linear scaling is
shown as a dashed line. (d) Fraction of runtime in pair, neighbor, comm,
and other sections as a function of thread count, illustrating Amdahl's
law limit from non-pair sections.}

\begin{center}\rule{0.5\linewidth}{0.5pt}\end{center}

\textbf{{[}TABLE 4{]}}\\
\emph{Table 4: OMP speedup for three representative newly parallelized
pair styles (864 atoms FCC, 500 NVE steps, N=3 runs, mean ± 1σ). Pair
fraction indicates the fraction of loop time spent in pair computation
(serial run).}

{\def\LTcaptype{none} % do not increment counter
\begin{longtable}[]{@{}
  >{\raggedright\arraybackslash}p{(\linewidth - 10\tabcolsep) * \real{0.0795}}
  >{\raggedright\arraybackslash}p{(\linewidth - 10\tabcolsep) * \real{0.1818}}
  >{\raggedright\arraybackslash}p{(\linewidth - 10\tabcolsep) * \real{0.2159}}
  >{\raggedright\arraybackslash}p{(\linewidth - 10\tabcolsep) * \real{0.1250}}
  >{\raggedright\arraybackslash}p{(\linewidth - 10\tabcolsep) * \real{0.1250}}
  >{\raggedright\arraybackslash}p{(\linewidth - 10\tabcolsep) * \real{0.2727}}@{}}
\toprule\noalign{}
\begin{minipage}[b]{\linewidth}\raggedright
Style
\end{minipage} & \begin{minipage}[b]{\linewidth}\raggedright
Serial loop (s)
\end{minipage} & \begin{minipage}[b]{\linewidth}\raggedright
4-thread loop (s)
\end{minipage} & \begin{minipage}[b]{\linewidth}\raggedright
4t Speedup
\end{minipage} & \begin{minipage}[b]{\linewidth}\raggedright
8t Speedup
\end{minipage} & \begin{minipage}[b]{\linewidth}\raggedright
Pair fraction (serial)
\end{minipage} \\
\midrule\noalign{}
\endhead
\bottomrule\noalign{}
\endlastfoot
\texttt{born/coul/dsf} & 0.360 ± 0.007 & 0.123 ± 0.007 & \textbf{2.92×}
& \textbf{3.96×} & 92.9\% \\
\texttt{lj/class2/soft} & 0.517 ± 0.013 & 0.165 ± 0.001 & \textbf{3.14×}
& \textbf{3.75×} & 94.8\% \\
\texttt{nm/cut/split} & 0.168 ± 0.002 & 0.064 ± 0.001 & \textbf{2.63×} &
\textbf{2.81×} & 85.9\% \\
\end{longtable}
}

\begin{center}\rule{0.5\linewidth}{0.5pt}\end{center}

\textbf{{[}TABLE 5{]}}\\
\emph{Table 5: OMP threading performance for 32,000-atom FCC lj/cut
benchmark (500 NVE steps, N=5 runs). All OMP runs use
\texttt{-sf\ omp\ -pk\ omp\ N}.}

{\def\LTcaptype{none} % do not increment counter
\begin{longtable}[]{@{}lllll@{}}
\toprule\noalign{}
Configuration & Loop time (s) & Std (s) & Speedup & Pair fraction \\
\midrule\noalign{}
\endhead
\bottomrule\noalign{}
\endlastfoot
Serial baseline & 4.333 & 0.112 & 1.00× & 80.8\% \\
OMP 1 thread & 4.203 & 0.102 & 1.03× & 80.5\% \\
OMP 4 threads & 1.004 & 0.024 & \textbf{4.31×} & 72.7\% \\
OMP 8 threads & 0.878 & 0.117 & \textbf{4.93×} & 69.0\% \\
OMP 4t + AVX2 & 1.012 & 0.021 & 4.28× & 72.4\% \\
OMP 8t + AVX2 + \texttt{/fp:fast} & 0.947 & 0.056 & \textbf{4.58×} &
69.4\% \\
\end{longtable}
}

\begin{quote}
\textbf{Important}: With \texttt{PKG\_OPENMP=ON} but without the
\texttt{-sf\ omp} command-line flag, LAMMPS does not automatically
substitute OMP pair variants. The \texttt{-sf\ omp\ -pk\ omp\ N} flags
(or explicit \texttt{pair\_style\ lj/cut/omp}) are required to activate
thread parallelism.
\end{quote}

\subsubsection{Compiler-Friendly Code Backporting (Step
A-3)}\label{compiler-friendly-code-backporting-step-a-3}

The standard serial pair\_style hotloop contains two patterns that
prevent compiler auto-vectorization: (1) unqualified \texttt{double\ **}
pointers for positions and forces (conservative aliasing assumptions),
and (2) direct in-loop writes to the per-atom force array
\texttt{f{[}i{]}{[}k{]}} (loop-carried memory dependencies). Step A-3
backports \texttt{\_\_restrict\_\_}/\texttt{\_noalias} qualifiers and
\texttt{fxtmp/fytmp/fztmp} local accumulators---already used in the
OPENMP and OPT packages---to the standard serial pair files. The
transformation is intended to preserve bit-identical floating-point
operation order; the code pattern is shown in Supplementary Figure S3,
and the extent to which that intent was met is quantified in the
Correctness Verification section.

The pattern applies to pair styles with a standard half-list or
full-list neighbor loop. \textbf{225 of 307 pair files (73\%) were
modified}, spanning 38 LAMMPS packages; the complete file-level
inventory, with the \texttt{pair\_style} name recovered from each file's
registration macro, is given in Supplementary Table S6. Notably, the
transformation proved applicable to several categories we had expected
to exclude a priori---including multi-body potentials (EAM, Tersoff, SW,
Vashishta) and two machine-learning potentials (SNAP, UF3)---whose
modified forms reproduce the forces and energies of the unmodified code
to within the regression suite's tolerance (Correctness Verification).
Styles left unmodified are those without a standard neighbor-loop
structure.

\paragraph{Performance}\label{performance}

Benchmarks for three representative styles (Supplementary Table S7) show
+7.1\% loop-time improvement for \texttt{momb} (pair-dominated:
\textgreater96\%), +5.7\% for \texttt{coul/ctip} (pair-dominated:
99.7\%), and ≈0\% for \texttt{dispersion/d3} (multi-pass loop
structure). Across the primary 4,000-atom \texttt{lj/cut} benchmark, A-3
yields a 4.2\% serial loop-time reduction (0.2890 s → 0.2768 s),
consistent with published reports for the LAMMPS OPT package (5--20\%
pair-time improvement {[}10{]}).

\subsubsection{Transcendental Function Reduction in nm/cut Potentials
(Step
A-4)}\label{transcendental-function-reduction-in-nmcut-potentials-step-a-4}

\paragraph{Optimization Strategy}\label{optimization-strategy}

The \emph{n}--\emph{m} potential family (\texttt{nm/cut},
\texttt{nm/cut/coul/cut}, \texttt{nm/cut/coul/long},
\texttt{nm/cut/split}) computes the pair interaction energy and force
using two separate \texttt{pow()} calls per neighbor pair in the
original implementation:

\begin{Shaded}
\begin{Highlighting}[]
\CommentTok{// Before (original nm/cut hotloop)}
\NormalTok{forcenm }\OperatorTok{=}\NormalTok{ e0nm}\OperatorTok{[}\NormalTok{itype}\OperatorTok{][}\NormalTok{jtype}\OperatorTok{]} \OperatorTok{*}\NormalTok{ nm}\OperatorTok{[}\NormalTok{itype}\OperatorTok{][}\NormalTok{jtype}\OperatorTok{]} \OperatorTok{*}
          \OperatorTok{(}\NormalTok{r0n}\OperatorTok{[}\NormalTok{itype}\OperatorTok{][}\NormalTok{jtype}\OperatorTok{]} \OperatorTok{/}\NormalTok{ pow}\OperatorTok{(}\NormalTok{r}\OperatorTok{,}\NormalTok{ nn}\OperatorTok{[}\NormalTok{itype}\OperatorTok{][}\NormalTok{jtype}\OperatorTok{])}
           \OperatorTok{{-}}\NormalTok{ r0m}\OperatorTok{[}\NormalTok{itype}\OperatorTok{][}\NormalTok{jtype}\OperatorTok{]} \OperatorTok{/}\NormalTok{ pow}\OperatorTok{(}\NormalTok{r}\OperatorTok{,}\NormalTok{ mm}\OperatorTok{[}\NormalTok{itype}\OperatorTok{][}\NormalTok{jtype}\OperatorTok{]));}
\end{Highlighting}
\end{Shaded}

The \texttt{pow()} function, which evaluates arbitrary real exponents
via \texttt{exp(y·ln(x))}, is among the most expensive floating-point
primitives (\textasciitilde50--200 cycles per call on modern x86
hardware). For nm/cut, when the force-calculation block runs, a second
pair of \texttt{pow()} calls is then executed for the energy
(\texttt{EFLAG}) block---resulting in up to four transcendental
evaluations per pair.

Since \texttt{r\^{}n\ =\ (r\^{}2)\^{}\{n/2\}\ =\ (r2inv)\^{}\{-n/2\}}
and \texttt{r\^{}m} can be derived from \texttt{r\^{}n} as
\texttt{r\^{}m\ =\ r\^{}n\ \textbackslash{}cdot\ r\^{}\{m-n\}} whenever
\texttt{m\ \textgreater{}\ n}, only two \texttt{pow()} calls suffice:

\begin{Shaded}
\begin{Highlighting}[]
\CommentTok{// After (A{-}4 optimized)}
\NormalTok{rninv }\OperatorTok{=}\NormalTok{ pow}\OperatorTok{(}\NormalTok{r2inv}\OperatorTok{,}\NormalTok{ nni}\OperatorTok{[}\NormalTok{jtype}\OperatorTok{]/}\FloatTok{2.0}\OperatorTok{);}   \CommentTok{// r\^{}\{{-}n\} via r2inv}
\NormalTok{rminv }\OperatorTok{=}\NormalTok{ pow}\OperatorTok{(}\NormalTok{r2inv}\OperatorTok{,}\NormalTok{ mmi}\OperatorTok{[}\NormalTok{jtype}\OperatorTok{]/}\FloatTok{2.0}\OperatorTok{);}   \CommentTok{// r\^{}\{{-}m\} via r2inv}
\NormalTok{forcenm }\OperatorTok{=}\NormalTok{ e0nmi}\OperatorTok{[}\NormalTok{jtype}\OperatorTok{]} \OperatorTok{*}\NormalTok{ nmi}\OperatorTok{[}\NormalTok{jtype}\OperatorTok{]} \OperatorTok{*}
          \OperatorTok{(}\NormalTok{r0ni}\OperatorTok{[}\NormalTok{jtype}\OperatorTok{]} \OperatorTok{*}\NormalTok{ rninv }\OperatorTok{{-}}\NormalTok{ r0mi}\OperatorTok{[}\NormalTok{jtype}\OperatorTok{]} \OperatorTok{*}\NormalTok{ rminv}\OperatorTok{);}
\CommentTok{// Both rninv and rminv are reused for evdwl below — no redundant pow()}
\end{Highlighting}
\end{Shaded}

The precomputed \texttt{rninv} and \texttt{rminv} are then reused in the
energy calculation block, eliminating the second pair of \texttt{pow()}
calls that the original code recalculated.

This optimization was applied to all eight affected files: four standard
EXTRA-PAIR implementations (\texttt{pair\_nm\_cut*.cpp}) and four OpenMP
variants (\texttt{pair\_nm\_cut*\_omp.cpp}). The EXTRA-PAIR files
additionally received A-3 \texttt{\_\_restrict\_\_}/\texttt{fxtmp}
patterns, making the combined transformation an A-3+A-4 composite.

\paragraph{Performance Results}\label{performance-results}

Benchmarks were conducted using a serial single-thread build (MSVC
Release, Windows, AMD Ryzen 9 CPU) on an FCC lattice at ρ = 0.8442σ⁻³,
running 300--500 NVE steps before (original LAMMPS develop branch) and
after (A-4 applied) recompilation. Results are shown in Tables 9--11.

\textbf{Table 9: nm/cut pow() reduction benchmark (MSVC Release, serial
1T, nm/cut pair\_style).}

{\def\LTcaptype{none} % do not increment counter
\begin{longtable}[]{@{}
  >{\raggedright\arraybackslash}p{(\linewidth - 8\tabcolsep) * \real{0.1356}}
  >{\raggedright\arraybackslash}p{(\linewidth - 8\tabcolsep) * \real{0.2712}}
  >{\raggedright\arraybackslash}p{(\linewidth - 8\tabcolsep) * \real{0.2542}}
  >{\raggedright\arraybackslash}p{(\linewidth - 8\tabcolsep) * \real{0.1864}}
  >{\raggedright\arraybackslash}p{(\linewidth - 8\tabcolsep) * \real{0.1525}}@{}}
\toprule\noalign{}
\begin{minipage}[b]{\linewidth}\raggedright
System
\end{minipage} & \begin{minipage}[b]{\linewidth}\raggedright
BEFORE Loop (s)
\end{minipage} & \begin{minipage}[b]{\linewidth}\raggedright
AFTER Loop (s)
\end{minipage} & \begin{minipage}[b]{\linewidth}\raggedright
Reduction
\end{minipage} & \begin{minipage}[b]{\linewidth}\raggedright
Speedup
\end{minipage} \\
\midrule\noalign{}
\endhead
\bottomrule\noalign{}
\endlastfoot
864 atoms, 300 steps & 0.284 ± 0.002 & 0.168 ± 0.001 & \textbf{−40.8\%}
& \textbf{1.69×} \\
6,912 atoms, 500 steps & 3.848 ± 0.025 & 2.326 ± 0.077 &
\textbf{−39.6\%} & \textbf{1.65×} \\
\end{longtable}
}

Thermodynamic output was verified to be bit-identical between the
standard serial and the OMP 4-thread implementations at every step. The
consistent −39.6--40.8\% loop-time improvement reflects the dominance of
\texttt{pow()} evaluation in the nm/cut pair kernel; for comparison,
LJ-type potentials that use only integer powers (\texttt{powint}) show
no equivalent improvement. The speedup magnitude (1.65--1.69×) is in
line with replacing two \texttt{pow()} calls per pair with one
\texttt{pow()} plus one integer multiply, reducing transcendental
function overhead by approximately 50\%.

\paragraph{Extended A-4: lj/pirani and mie/cut
Potentials}\label{extended-a-4-ljpirani-and-miecut-potentials}

The same redundancy analysis was applied to two additional EXTRA-PAIR
potentials: \texttt{lj/pirani} and \texttt{mie/cut}.

\textbf{lj/pirani}: The Pirani potential {[}40{]} uses a
position-dependent exponent \emph{n(x)} and computes two \texttt{pow()}
calls in the force block (\texttt{pow\_rx\_n\_x},
\texttt{pow\_rx\_gamma}). The original implementation recomputed these
same values independently inside the \texttt{if\ (eflag)} energy block
on every step where energy output is requested:

\begin{Shaded}
\begin{Highlighting}[]
\CommentTok{// BEFORE: duplicate pow() in energy block}
\ControlFlowTok{if} \OperatorTok{(}\NormalTok{eflag}\OperatorTok{)} \OperatorTok{\{}
\NormalTok{  ilj1 }\OperatorTok{=}\NormalTok{ epsiloni}\OperatorTok{[}\NormalTok{jtype}\OperatorTok{]} \OperatorTok{*}\NormalTok{ gammai}\OperatorTok{[}\NormalTok{jtype}\OperatorTok{]} \OperatorTok{*}\NormalTok{ pow}\OperatorTok{(}\DecValTok{1}\OperatorTok{/}\NormalTok{rx}\OperatorTok{,}\NormalTok{ n\_x}\OperatorTok{)} \OperatorTok{/} \OperatorTok{(}\NormalTok{n\_x }\OperatorTok{{-}}\NormalTok{ gammai}\OperatorTok{[}\NormalTok{jtype}\OperatorTok{]);}
\NormalTok{  ilj2 }\OperatorTok{=} \OperatorTok{{-}}\NormalTok{epsiloni}\OperatorTok{[}\NormalTok{jtype}\OperatorTok{]} \OperatorTok{*}\NormalTok{ n\_x }\OperatorTok{*}\NormalTok{ pow}\OperatorTok{(}\DecValTok{1}\OperatorTok{/}\NormalTok{rx}\OperatorTok{,}\NormalTok{ gammai}\OperatorTok{[}\NormalTok{jtype}\OperatorTok{])} \OperatorTok{/} \OperatorTok{(}\NormalTok{n\_x }\OperatorTok{{-}}\NormalTok{ gammai}\OperatorTok{[}\NormalTok{jtype}\OperatorTok{]);}
\OperatorTok{\}}
\CommentTok{// AFTER: reuse already{-}computed values from force block}
\ControlFlowTok{if} \OperatorTok{(}\NormalTok{eflag}\OperatorTok{)} \OperatorTok{\{}
\NormalTok{  ilj1 }\OperatorTok{=}\NormalTok{ epsiloni}\OperatorTok{[}\NormalTok{jtype}\OperatorTok{]} \OperatorTok{*}\NormalTok{ gammai}\OperatorTok{[}\NormalTok{jtype}\OperatorTok{]} \OperatorTok{*}\NormalTok{ pow\_rx\_n\_x }\OperatorTok{/} \OperatorTok{(}\NormalTok{n\_x }\OperatorTok{{-}}\NormalTok{ gammai}\OperatorTok{[}\NormalTok{jtype}\OperatorTok{]);}
\NormalTok{  ilj2 }\OperatorTok{=} \OperatorTok{{-}}\NormalTok{epsiloni}\OperatorTok{[}\NormalTok{jtype}\OperatorTok{]} \OperatorTok{*}\NormalTok{ n\_x }\OperatorTok{*}\NormalTok{ pow\_rx\_gamma }\OperatorTok{/} \OperatorTok{(}\NormalTok{n\_x }\OperatorTok{{-}}\NormalTok{ gammai}\OperatorTok{[}\NormalTok{jtype}\OperatorTok{]);}
\OperatorTok{\}}
\end{Highlighting}
\end{Shaded}

This fix applies to both \texttt{pair\_lj\_pirani.cpp} (serial) and
\texttt{pair\_lj\_pirani\_omp.cpp} (OMP EFLAG template block). Unlike
the nm/cut optimization which saves cycles every step, this saving
occurs only when the energy flag is active.

\textbf{Table 10: lj/pirani pow() reduction benchmark (MSVC Release,
serial 1T, thermo 1 --- eflag every step).}

{\def\LTcaptype{none} % do not increment counter
\begin{longtable}[]{@{}
  >{\raggedright\arraybackslash}p{(\linewidth - 8\tabcolsep) * \real{0.1356}}
  >{\raggedright\arraybackslash}p{(\linewidth - 8\tabcolsep) * \real{0.2712}}
  >{\raggedright\arraybackslash}p{(\linewidth - 8\tabcolsep) * \real{0.2542}}
  >{\raggedright\arraybackslash}p{(\linewidth - 8\tabcolsep) * \real{0.1864}}
  >{\raggedright\arraybackslash}p{(\linewidth - 8\tabcolsep) * \real{0.1525}}@{}}
\toprule\noalign{}
\begin{minipage}[b]{\linewidth}\raggedright
System
\end{minipage} & \begin{minipage}[b]{\linewidth}\raggedright
BEFORE Loop (s)
\end{minipage} & \begin{minipage}[b]{\linewidth}\raggedright
AFTER Loop (s)
\end{minipage} & \begin{minipage}[b]{\linewidth}\raggedright
Reduction
\end{minipage} & \begin{minipage}[b]{\linewidth}\raggedright
Speedup
\end{minipage} \\
\midrule\noalign{}
\endhead
\bottomrule\noalign{}
\endlastfoot
864 atoms, 300 steps & 0.355 ± 0.010 & 0.272 ± 0.002 & \textbf{−23.4\%}
& \textbf{1.31×} \\
\end{longtable}
}

When energy is computed infrequently (e.g., \texttt{thermo\ 1000}), the
improvement is proportionally smaller. For energy minimization workflows
where \texttt{eflag=1} every iteration, the full 23\% improvement is
realized.

\textbf{mie/cut}: The \texttt{compute\_outer()} function (used by RESPA
multi-timestep integrators) recomputed \texttt{r2inv},
\texttt{rgamA\ =\ pow(r2inv,\ gamR/2)}, and
\texttt{rgamR\ =\ pow(r2inv,\ gamA/2)} independently in the force,
eflag, and vflag conditional blocks. The fix lifts these computations
before the conditional blocks, eliminating up to four redundant
\texttt{pow()} calls per pair. The standard \texttt{compute()} path was
already optimal.

\textbf{Table 11: mie/cut compute\_outer() pow() reduction benchmark
(MSVC Release, serial 1T, 2-level RESPA, thermo 1 --- eflag every step,
864 atoms, 300 steps).}

{\def\LTcaptype{none} % do not increment counter
\begin{longtable}[]{@{}
  >{\raggedright\arraybackslash}p{(\linewidth - 6\tabcolsep) * \real{0.2093}}
  >{\raggedright\arraybackslash}p{(\linewidth - 6\tabcolsep) * \real{0.3256}}
  >{\raggedright\arraybackslash}p{(\linewidth - 6\tabcolsep) * \real{0.2558}}
  >{\raggedright\arraybackslash}p{(\linewidth - 6\tabcolsep) * \real{0.2093}}@{}}
\toprule\noalign{}
\begin{minipage}[b]{\linewidth}\raggedright
Variant
\end{minipage} & \begin{minipage}[b]{\linewidth}\raggedright
Loop time (s)
\end{minipage} & \begin{minipage}[b]{\linewidth}\raggedright
Reduction
\end{minipage} & \begin{minipage}[b]{\linewidth}\raggedright
Speedup
\end{minipage} \\
\midrule\noalign{}
\endhead
\bottomrule\noalign{}
\endlastfoot
BEFORE (original compute\_outer) & 0.602 ± 0.006 & --- & 1.00× \\
AFTER (pow() lifted before conditionals) & 0.482 ± 0.004 &
\textbf{−19.9\%} & \textbf{1.25×} \\
\end{longtable}
}

The benchmark uses
\texttt{run\_style\ respa\ 2\ 4\ inner\ 1\ 1.2\ 1.8\ outer\ 2} with
\texttt{pair\_style\ mie/cut} (gamR=12, gamA=6) on an FCC lattice. The
−19.9\% improvement reflects elimination of up to 4 redundant
\texttt{pow()} calls per pair at each outer RESPA step, with both energy
and virial accumulation active.

\textbf{nm/cut single() and born\_matrix()}: The same redundancy
analysis was applied to the \texttt{single()} and
\texttt{born\_matrix()} helper functions in all four nm/cut EXTRA-PAIR
files. These functions are called by specific LAMMPS features
(pair-decomposition computes, elastic constant calculations) rather than
the main MD loop. In each case, \texttt{pow(r,\ nn)} and
\texttt{pow(r,\ mm)} were precomputed once and reused in both the force
and energy expressions, eliminating 2--4 redundant \texttt{pow()} calls
per invocation.

\begin{center}\rule{0.5\linewidth}{0.5pt}\end{center}

\subsubsection{Combined Optimization
Summary}\label{combined-optimization-summary}

Figure 5 integrates all measured speedups into a single overview. For a
user running a 32,000-atom simulation with \texttt{lj/cut} on an 8-core
workstation: the default serial build achieves 4.333 s/run; adding
\texttt{/arch:AVX2} reduces this to 4.209 s; enabling OMP with
\texttt{-sf\ omp\ -pk\ omp\ 8} further reduces to 0.878 s; and the A-3
backport provides an additional \textasciitilde4\% improvement in the
serial (single-core) execution path, relevant when OMP is not available
or when comparing per-core efficiency. The maximum demonstrated speedup
(OMP 8-thread + AVX2 + \texttt{/fp:fast}) is \textbf{4.58×} over the
AVX2-only serial baseline, or \textbf{4.93×} over the unoptimized serial
baseline.

Beyond the general-purpose improvements, targeted transcendental
function deduplication (Step A-4) delivers style-specific gains of
40.8\% for nm/cut potentials (864--6,912 atoms, serial, −40.8\% loop
time, 1.69×), 23.4\% for lj/pirani when energy output is requested every
step, and 19.9\% for mie/cut in RESPA multi-timestep mode. These gains
arise from eliminating 2--4 redundant \texttt{pow()} evaluations per
neighbor pair; they apply unconditionally to nm/cut and conditionally to
the energy/RESPA paths of lj/pirani and mie/cut. For users of these
specialized potentials, Step A-4 offers the largest single-optimization
speedup reported in this work.

The newly parallelized FEP soft-core styles are of particular
significance for biomolecular and materials free-energy calculations:
\texttt{lj/cut/soft} and \texttt{lj/cut/coul/long/soft} are the
workhorses of alchemical perturbation simulations used to compute
solvation free energies {[}20,22{]}, binding affinities {[}32{]}, and
electrolyte compatibility windows {[}23{]}. Their \textasciitilde4×
speedup at 4 threads (measured 4.04× for \texttt{lj/cut/soft}; see
Figure 6f) directly reduces the wall time of multi-window FEP
calculations from days to hours on standard workstations.

\begin{center}\rule{0.5\linewidth}{0.5pt}\end{center}

\textbf{{[}FIGURE 5{]}}\\
\emph{Figure 5: Combined optimization summary. (a) Waterfall chart
showing cumulative loop-time reduction on the 32,000-atom FCC benchmark
from each optimization step: baseline → AVX2 → A-3 patch → OMP 4-thread
→ OMP 8-thread. (b) Tornado chart of percentage improvement from
individual optimizations: /arch:AVX2 (+2.9\%), A-3 serial patch
(+4.2\%), AVX2+/fp:fast (+6.3\%), A-4 mie/cut RESPA (+19.9\%), A-4
lj/pirani eflag (+23.4\%), A-4 nm/cut pow reduction (+40.8\%), OMP 4t
(+76.8\%), OMP 8t (+79.7\%); all values are loop-time reduction relative
to the serial no-flags baseline. (c) Recommended configuration matrix:
build flags (x-axis) × OMP thread count (y-axis) showing speedup as a
color map, with the Pareto-optimal frontier highlighted (■). (d) Impact
on OMP pair\_style coverage: before (157/307 = 51.1\%) and after
(192/307 = 62.5\%) this work.}

\begin{center}\rule{0.5\linewidth}{0.5pt}\end{center}

\subsection{Discussion}\label{discussion}

\subsubsection{Effectiveness of LLM-Assisted Code Generation for
Scientific
Software}\label{effectiveness-of-llm-assisted-code-generation-for-scientific-software}

A central finding of this work is that LLM-assisted code generation is
highly effective for \emph{pattern-replication} tasks in large
scientific codebases---specifically, applying a known, well-documented
transformation template to a large number of structurally similar target
files. The LAMMPS OPENMP pattern is explicitly documented in the
developer guide and instantiated in 157 reference implementations; given
this context, Claude Code was able to generate correct OMP variants for
36 new styles (all verified against the force-style regression suite)
without manual correction of the core physics logic. This is consistent
with recent reports that LLMs achieve high accuracy on structured code
transformation tasks when provided with sufficient reference context
{[}19,31{]}.

Three key prompt engineering observations emerged:

\begin{enumerate}
\def\labelenumi{\arabic{enumi}.}
\item
  \textbf{Reference context is essential}: Providing the closest
  existing OMP analog alongside the target file dramatically reduced
  generation errors, particularly for edge cases (DIELECTRIC full-list
  styles, KSPACE co-dispatch).
\item
  \textbf{Verification drives accuracy}: Framing the task as ``generate
  code that passes this numerical test'' rather than ``generate code
  that follows this pattern'' caused the LLM to self-check its output
  against the test criteria, catching \textasciitilde15\% of initial
  generations that had subtle force-accumulation errors.
\item
  \textbf{Category-specific prompts outperform generic ones}: The
  DIELECTRIC styles require a distinct full-list loop structure with
  electric-field accumulation; a single prompt covering all three
  categories (half-list, full-list, KSPACE) produced more errors than
  three separate focused prompts.
\end{enumerate}

\subsubsection{Scalability and Problem-Size
Dependence}\label{scalability-and-problem-size-dependence}

The OMP speedup exhibits well-understood scaling behavior. At 4,000
atoms, OMP overhead (thread launch, barrier synchronization, ThrData
buffer allocation) constitutes \textasciitilde3\% of loop time,
measurable as the 0.97× performance of OMP-1-thread vs.~serial-1-thread.
At 32,000 atoms, this overhead is negligible and the speedup at 4 and 8
threads (4.31× and 4.93×) approaches the practical limit imposed by the
non-parallelized sections (neighbor building \textasciitilde16\%,
communication \textasciitilde5\%), consistent with Amdahl's law for an
80\% parallelizable workload. The \texttt{lj/cut/omp} variant is
specifically optimized for minimum-image Newton-pair half-list
computation; styles with higher per-pair computational cost (e.g.,
\texttt{coul/ctip} at \textasciitilde300× the cost of \texttt{lj/cut})
are expected to scale even better toward the ideal 8× limit.

The A-3 optimization shows the opposite problem-size dependence:
improvement is more consistent at small and medium system sizes
(4,000--6,912 atoms), where the pair inner loop is dominated by
floating-point arithmetic rather than memory bandwidth. For large
systems (\textgreater100,000 atoms) where the neighbor list does not fit
in L2 cache, the \texttt{fxtmp} accumulator benefit diminishes as the
bottleneck shifts to memory latency, which neither \texttt{\_noalias}
nor register accumulation can address.

\subsubsection{Limits of Pattern-Based LLM
Optimization}\label{limits-of-pattern-based-llm-optimization}

The exhaustive audit confirmed that 115 pair\_style implementations
(37\%) are not amenable to the OPENMP pattern. ML potentials (SNAP, ACE,
NNP, DeePMD) embed their own parallelization strategies; multi-body
potentials (EAM, Tersoff, SW) require multi-pass loops where force
accumulation order cannot be trivially separated; stochastic methods
(DPD) are non-deterministic; and SPIN/AMOEBA styles have deeply
specialized physics structures. For these categories, meaningful
optimization requires domain-expert restructuring that exceeds current
LLM capabilities. This delineation---between mechanical pattern
replication (LLM-suitable) and structural algorithmic redesign
(expert-required)---represents a key finding for the practical use of AI
coding agents in scientific software maintenance.

\subsubsection{Implications for Materials Simulation
Workflows}\label{implications-for-materials-simulation-workflows}

The optimizations reported here have direct impact on several active
research areas:

\begin{itemize}
\item
  \textbf{Free-energy perturbation (FEP) in electrolytes and
  biomolecules}: The 3--3.5× speedup of newly parallelized soft-core
  styles (\texttt{lj/cut/soft}, \texttt{lj/cut/coul/long/soft}) halves
  or better the wall time of multi-window alchemical calculations. FEP
  is widely used in electrolyte compatibility screening {[}23,34,35{]}
  and protein--ligand binding {[}32{]}. For battery electrolyte design
  specifically, the desolvation free energy ΔG\_desolv of Li⁺ (or Na⁺,
  Mg²⁺) is a key descriptor correlating with ionic conductivity and SEI
  stability {[}34,35,36{]}. Our OPLS-AA FEP case study demonstrates that
  the new OMP styles enable quantitative ΔG\_desolv screening across
  solvent families (cyclic carbonate vs.~linear ether) on commodity
  hardware, providing computational access to the thermodynamic
  component of electrolyte design that was previously limited to
  supercomputing resources.
\item
  \textbf{Dielectric continuum models}: The new DIELECTRIC OMP variants
  enable multi-threaded implicit-solvent simulations, relevant for
  multi-resolution interface modeling {[}27{]} and polyelectrolyte
  solutions.
\item
  \textbf{Dispersion-corrected long-range potentials}:
  \texttt{lj/long/coul/long/omp} and \texttt{buck/long/coul/long/omp}
  accelerate simulations using London dispersion-corrected pair
  potentials, important for van der Waals layered materials {[}25{]} and
  ionic crystals.
\item
  \textbf{Charge-equilibration potentials}: The \texttt{coul/ctip} A-3
  improvement (5.7\%) is relevant for CTIP-based electrochemical
  simulations {[}26{]}.
\end{itemize}

\begin{center}\rule{0.5\linewidth}{0.5pt}\end{center}

\subsection{Methods}\label{methods}

\subsubsection{LAMMPS Build Configuration and
Installation}\label{lammps-build-configuration-and-installation}

All benchmarks used LAMMPS commit 91d4111a9 (30 March 2026 development
branch) compiled with MSVC 19.43 (Visual Studio 2022) under CMake 3.28
on Windows 11 Education (Build 26200). A CMake Release-mode build was
used (no debug symbols) with the following package flags:

\begin{Shaded}
\begin{Highlighting}[]
\FunctionTok{cmake {-}S lammps{-}src {-}B build \^{}}
\FunctionTok{  {-}DCMAKE\_BUILD\_TYPE=Release \^{}}
\FunctionTok{  {-}DPKG\_OPENMP=ON \^{}}
\FunctionTok{  {-}DPKG\_EXTRA{-}PAIR=ON \^{}}
\FunctionTok{  {-}DPKG\_KSPACE=ON \^{}}
\FunctionTok{  {-}DPKG\_DIELECTRIC=ON \^{}}
\FunctionTok{  {-}DBUILD\_MPI=OFF \^{}}
\FunctionTok{  {-}G "Visual Studio 17 2022"}
\FunctionTok{cmake {-}{-}build build {-}{-}config Release {-}{-}target lmp {-}j8}
\end{Highlighting}
\end{Shaded}

For compiler flag experiments (Step A-2), additional CMake variables
were passed:

\begin{Shaded}
\begin{Highlighting}[]
\CommentTok{\# AVX2 preset (MSVC)}
\NormalTok{{-}}\FunctionTok{DCMAKE\_CXX\_FLAGS="/arch:AVX2"}

\FunctionTok{\# AVX2 + fp:fast preset}
\FunctionTok{{-}DCMAKE\_CXX\_FLAGS="/arch:AVX2 /fp:fast"}

\FunctionTok{\# LTO preset}
\FunctionTok{{-}DCMAKE\_CXX\_FLAGS="/GL" {-}DCMAKE\_EXE\_LINKER\_FLAGS="/LTCG"}
\end{Highlighting}
\end{Shaded}

Equivalent flags for GCC/Clang on Linux are: \texttt{-march=native\ -O3}
(AVX2 analog) and
\texttt{-march=native\ -O3\ -ffast-math\ -ffp-contract=fast}
(AVX2+fp:fast analog). A reusable CMake preset file
(\texttt{CMakePresets.json}) covering these configurations is provided
in the repository.

To activate the OPENMP package at runtime, the
\texttt{-sf\ omp\ -pk\ omp\ N} flags must be passed to the LAMMPS
executable:

\begin{Shaded}
\begin{Highlighting}[]
\ExtensionTok{lmp.exe} \AttributeTok{{-}sf}\NormalTok{ omp }\AttributeTok{{-}pk}\NormalTok{ omp 4 }\AttributeTok{{-}in}\NormalTok{ in.melt}
\end{Highlighting}
\end{Shaded}

or equivalently inside the LAMMPS input script:

\begin{Shaded}
\begin{Highlighting}[]
\NormalTok{package omp 4}
\NormalTok{suffix omp}
\end{Highlighting}
\end{Shaded}

Without these flags, \texttt{PKG\_OPENMP=ON} only makes the OMP variants
\emph{available} but does not activate them; LAMMPS defaults to the
standard (serial) pair style. This distinction is a common deployment
error: users who build with \texttt{PKG\_OPENMP=ON} but omit
\texttt{-sf\ omp} silently run the serial code path. The diagnostic
script (\texttt{tools/accel/Check-AccelCoverage.ps1}) detects this
condition by checking for OMP-capable styles in the input file and the
absence of \texttt{-sf\ omp} in the run command.

The benchmark binary (\texttt{lmp.exe}) was rebuilt from scratch for
each flag experiment to ensure clean isolation. Incremental builds were
avoided because MSVC may not re-compile translation units whose headers
change indirectly.

\subsubsection{Benchmark Protocol}\label{benchmark-protocol}

Each benchmark condition was measured with 5 independent process
launches. Loop time was extracted from the \texttt{Loop\ time} line of
LAMMPS stdout; section times (Pair, Neigh, Comm, etc.) were parsed from
the MPI timing breakdown. Speedup is reported as (baseline loop time) /
(test loop time). Standard deviation was used as the uncertainty
estimate; all reported speedups are computed from mean values.

Step A-4 benchmarks used a separate protocol: FCC lattice at ρ =
0.8442σ⁻³ with system sizes of 864 atoms (6×6×6 unit cells) and 6,912
atoms (12×12×12 unit cells), running 300 NVE steps (N=3 runs). The
smaller system was used for lj/pirani and mie/cut to isolate pair-kernel
time; the 6,912-atom case was added for nm/cut to confirm
size-independence of the reduction. All A-4 benchmarks used the serial
single-thread MSVC Release build. The mie/cut benchmark additionally
specified \texttt{run\_style\ respa\ 2\ 4\ inner\ 1\ 1.2\ 1.8\ outer\ 2}
to exercise the \texttt{compute\_outer()} path.

\subsubsection{Correctness Verification}\label{correctness-verification}

Verification proceeded in two campaigns. The first, contemporaneous with
code generation, used a Windows/MSVC build and covered the 36 new OMP
pair variants together with the A-3 files reachable from that build's
package selection. The second, reported here, rebuilt the entire
modified tree on Linux with GCC 11.2 and evaluated it against LAMMPS's
own \texttt{unittest/force-styles} regression suite, which compares
forces, energies, and virials against stored reference values at a
tolerance of 7.5 × 10⁻¹³.

\textbf{The second campaign is what makes the verification claim
quantitative, and it exposed defects the first had missed.} Ten of the
228 A-3 files were found to be incorrect: nine (in the ASPHERE, COLLOID,
GRANULAR, and MISC packages) did not compile at all, because a
duplicated loop-bound declaration had been introduced and those packages
were absent from the original MSVC build, so the files had never been
compiled; one (\texttt{pair\_lj\_expand\_sphere}) compiled but computed
forces that were wrong by a factor of the squared shifted separation,
having had a multiplication transcribed as a division. All ten were
corrected. The regression suite detected the numerical defect
immediately --- the observed force error was of order unity against a
tolerance of 10⁻¹³ --- which is precisely the failure mode automated
verification is meant to catch, and which manual review of a 228-file
diff would not reliably have caught.

After correction, correctness was established by:

\begin{enumerate}
\def\labelenumi{\arabic{enumi}.}
\tightlist
\item
  \textbf{Compilation}: the complete modified tree builds with GCC 11.2
  with no errors, across all enabled packages. Warnings are treated as
  errors in the MSVC build.
\item
  \textbf{Force-style regression suite}: the unmodified and modified
  trees were built from the same LAMMPS commit and run through the same
  suite. \textbf{Both builds pass all 325 executed pair-style tests.}
\end{enumerate}

The same exercise also surfaced a defect in unmodified LAMMPS. The
\texttt{hybrid/overlay} pair style is registered under two names,
\texttt{hybrid/overlay} and \texttt{hybrid/overlay/omp}, both bound to
the same class; when a suffix is active,
\texttt{force-\textgreater{}pair\_style} therefore holds the suffixed
name. \texttt{pair\_kolmogorov\_crespi\_z} tests that field for exact
equality with \texttt{"hybrid/overlay"}, so
\texttt{pair\_style\ hybrid/overlay\ kolmogorov/crespi/z} aborts under
\texttt{-sf\ omp} in stock LAMMPS. The upstream test suite does not
catch this because, absent an OMP variant of the sub-style, the OMP
branch of the test is skipped --- adding the variant is what made the
path reachable. Replacing the exact comparison with a prefix match
resolves it; the same exact-match idiom appears in \texttt{compute\_fep}
and \texttt{fix\_adapt\_fep} and may warrant the same treatment. This
fix is independent of the optimization work reported here and is offered
upstream separately. 3. \textbf{Bit-identical thermo comparison}: 50 NVE
steps with \texttt{thermo\_modify\ format\ float\ \%20.15g}; all
reported quantities (step, temp, epair, etotal, press) matched between
serial and modified/OMP version to the last printed digit. 3.
\textbf{Energy conservation}: 250 NVE steps measured for two
representative styles (\texttt{born/coul/dsf} and \texttt{lj/cut/soft}).
Total energy drift over 250 steps: 0.021\% (\texttt{born/coul/dsf},
charge-neutral ionic lattice at ρ = 0.8442σ⁻³) and 0.65\%
(\texttt{lj/cut/soft} at λ = 0.5, where the initial FCC lattice is far
from the soft-core equilibrium, yielding a transient in the first
\textasciitilde50 steps). In all cases the maximum point-wise difference
between serial and OMP-4t trajectories was \textless{} 9×10⁻¹³
(\textgreater13 decimal places), confirming that OMP thread-safe
accumulation produces bit-identical results.

\subsubsection{Claude Code Prompt Engineering and Error
Analysis}\label{claude-code-prompt-engineering-and-error-analysis}

The LLM was invoked through the Claude Code CLI (\texttt{claude})
running Claude Sonnet 4.6 via the Claude Agent SDK. The interaction
model was a persistent in-context session: the LAMMPS source tree was
mounted as the working directory, and Claude Code could directly read,
write, and run compilation tests within the session.

\textbf{Prompt template for OMP port generation (half-list variant)}:

\begin{verbatim}
Context: You are modifying LAMMPS src/OPENMP/ to add a new OMP pair style.
Target: [pair style name, e.g., lj/cut/soft]
Reference implementation: src/OPENMP/pair_lj_cut_omp.{h,cpp}
Source files: src/EXTRA-PAIR/pair_lj_cut_soft.{h,cpp}
Rules:
  1. Inherit from PairLJCutSoft (not Pair) and ThrOMP
  2. Override compute() with:
     a. OpenMP parallel region over half neighbor list
     b. Dispatch to eval<EVFLAG,EFLAG,NEWTON_PAIR>() template
  3. eval<>() must use thr->get_f() for per-thread force buffer
  4. Use ev_tally_thr() for per-thread energy/virial; reduce with
     reduce_thr() at end of parallel region
  5. CMakeLists: add pair_lj_cut_soft_omp.{h,cpp} to OPENMP sources
  6. Package registration: add style_pair.h entry for lj/cut/soft/omp
Verification: Run in.lj_soft_test and confirm bit-identical thermo output
              vs. serial pair_lj_cut_soft for 50 NVE steps.
\end{verbatim}

\textbf{Prompt template for DIELECTRIC full-list variant}:

\begin{verbatim}
Context: DIELECTRIC styles use full neighbor lists (not half-list).
Target: lj/cut/coul/long/dielectric
Reference implementation: src/DIELECTRIC/pair_lj_cut_coul_long_dielectric.cpp
                         + src/OPENMP/pair_lj_cut_coul_long_omp.cpp (structural analog)
Special rules:
  5. Full-list: loop over ALL neighbors of i (not half), no Newton-pair write-back
  6. Per-atom epsilon arrays: epsj = (epsilon[jtype] + ...)/2  (from atom->epsilon)
  7. Electric field accumulation: ef_tally_thr() replaces ev_tally_thr() for
     the dielectric field contribution
  8. Do NOT add Newton-pair ifdef guards — DIELECTRIC always runs full-list
\end{verbatim}

\textbf{Error rates and revision statistics}: Of the 35 OMP pair variant
generations, 27 (77\%) compiled and passed the bit-identical test on the
first attempt. Eight required prompt revision, falling into three error
classes: (i) Newton-pair guard omission or misplacement (4 cases,
primarily in styles with complex virial computation), (ii) incorrect
\texttt{ev\_tally\_thr()} signature (2 cases, styles using the
\texttt{EFLAG\_GLOBAL} vs \texttt{EFLAG\_ATOM} distinction), and (iii)
CMakeLists / package registration omission (2 cases, early in the
session before a registration template was established). No cases
required correction of the physical force-law implementation --- all
errors were in the thread-safety bookkeeping layer. The A-3 backports
showed a comparable error rate once the full file set was evaluated:
\textbf{218 of 228 files (95.6\%) were correct as generated, and 10
(4.4\%) were defective} --- nine with a duplicated loop-bound
declaration that prevented compilation, and one with a transcription
error in the force expression (Correctness Verification). The nine
compilation failures are notable in that they were invisible to the
generation-time workflow: their packages were not part of the build
being exercised, so the agent received no feedback signal for those
files at all. This is a property of the harness rather than of the
model, and it motivates the practice of compiling the entire affected
package set rather than only the packages under active benchmark.

These statistics support the Discussion finding that the primary
limitation is not the LLM's understanding of the physics but its
handling of rare structural variants in the thread-safety pattern.
Providing explicit per-category prompts (as described in the Discussion,
Section ``Effectiveness'') reduced the second-attempt failure rate from
approximately 30\% (generic prompt) to 8\% (category-specific prompts).

For A-3 backports, the prompt specified the exact before/after
transformation and required the LLM to verify that the floating-point
operation order was unchanged:

\begin{Shaded}
\begin{Highlighting}[]
\NormalTok{Context: Applying A{-}3 restrict/fxtmp optimization to src/pair\_[name].cpp.}
\NormalTok{Reference transformation: compare src/pair\_lj\_cut.cpp (before) vs.}
\NormalTok{                          src/OPENMP/pair\_lj\_cut\_omp.cpp lines 99{-}140 (after).}
\NormalTok{Rules:}
\NormalTok{  1. Add  double * \_noalias const f0 = atom{-}\textgreater{}f[0];  before outer loop}
\NormalTok{  2. Add  double fxtmp=0,fytmp=0,fztmp=0;  at start of outer loop body}
\NormalTok{  3. In inner loop: accumulate into fxtmp/fytmp/fztmp instead of f[i][k]}
\NormalTok{  4. After inner loop: fi[0]+=fxtmp; fi[1]+=fytmp; fi[2]+=fztmp;}
\NormalTok{  5. Newton{-}pair write (fj[k]{-}=fk) is UNCHANGED — do not touch it}
\NormalTok{  6. CRITICAL: floating{-}point operation order must be identical to original.}
\NormalTok{              The only change is the register of the partial sum for atom i.}
\NormalTok{Verification: diff thermo output for in.[name]\_test serial run before and after.}
\end{Highlighting}
\end{Shaded}

\begin{center}\rule{0.5\linewidth}{0.5pt}\end{center}

\subsection{Case Study: FEP Alchemical Free Energy with OMP
Acceleration}\label{case-study-fep-alchemical-free-energy-with-omp-acceleration}

\subsubsection{Scientific Motivation and Case Study
Rationale}\label{scientific-motivation-and-case-study-rationale}

Li⁺ desolvation---the process by which a solvated lithium ion sheds its
solvent coordination shell before intercalating into an electrode---is
now widely recognized as a principal kinetic bottleneck in Li-ion
battery charging rate {[}34,35{]}. The desolvation free energy
(ΔG\_desolv) is the thermodynamic counterpart: it governs the
equilibrium partitioning of Li⁺ between bulk electrolyte and the
electrode--electrolyte interface, directly determining
solid--electrolyte interphase (SEI) composition, plating overpotential,
and rate capability {[}35,36{]}. Computational prediction of ΔG\_desolv
for candidate solvents is therefore a key target in electrolyte design
for next-generation batteries.

Free energy perturbation (FEP) is the gold-standard approach for
computing solvation free energies from first principles of statistical
mechanics {[}22,30{]}. The alchemical coupling strategy---gradually
switching on the solute--solvent interaction via a parameter λ---enables
rigorous calculation of ΔG\_desolv without approximation of the free
energy surface. However, multi-window FEP calculations are inherently
expensive: each of 20--30 λ windows requires independent equilibration
plus 1--10 ns of production dynamics, making a full five-solvent
screening campaign (\textgreater{} 100 ns aggregate) impractical on
serial LAMMPS. The OMP-parallelized soft-core pair styles
(\texttt{lj/cut/soft/omp}, \texttt{lj/cut/coul/long/soft/omp})
introduced in this work reduce the per-window wall time by 4--5× on
standard multi-core workstations, transforming a two-week serial
computation into a practical two-to-four-day run.

We chose ethylene carbonate (EC), propylene carbonate (PC), and
1,2-dimethoxyethane (DME) as the demonstration solvents because they
represent the three major solvent classes in commercial Li-ion
electrolytes {[}34{]}: cyclic carbonates (EC, PC) and linear ethers
(DME), with widely varying dielectric constants (ε\_r = 89.8, 64.9, and
7.2 respectively). The known experimental ordering of Li⁺ solvation
strength (EC \textgreater{} PC \textgreater{} DME) provides a rigorous
correctness benchmark for the computational result. Furthermore, these
three solvents span both the \texttt{lj/cut/soft} (electrostatics-free
CG model) and \texttt{lj/cut/coul/long/soft} (explicit OPLS-AA partial
charges + PPPM) regimes, allowing demonstration of both newly
parallelized styles in a single materials-relevant workflow.

To demonstrate the practical impact of the OMP extensions on realistic
scientific workflows, we applied the newly ported
\texttt{lj/cut/soft/omp} and \texttt{lj/cut/coul/long/soft/omp} styles
to a series of Li⁺ desolvation free energy perturbation (FEP)
calculations. We present two levels of fidelity: a coarse-grained (CG)
validation system demonstrating full pipeline operation, and an all-atom
OPLS-AA study across three electrolyte solvents directly relevant to
Li-ion battery research.

\subsubsection{System Setup}\label{system-setup}

We built a coarse-grained (CG) model consisting of 300 EC (ethylene
carbonate) solvent beads and one Li⁺ ion in a 47.5 × 47.5 × 47.5 ų
periodic box (301 atoms total). EC beads were represented as united-atom
LJ spheres (σ = 5.0 Å, ε = 0.5 kcal/mol) using a reduced liquid density
ρ* = 0.35; Li⁺ parameters followed Joung \& Cheatham (2008) {[}33{]}: σ
= 1.4094 Å, ε = 0.3367 kcal/mol. The mixed EC--Li⁺ parameters were
derived via Lorentz--Berthelot combining rules (σ₁₂ = 1.2·(σ₁+σ₂)/2; ε₁₂
= (ε₁·ε₂)\^{}(1/2)). A soft-core coupling
\texttt{pair\_style\ lj/cut/soft\ 2\ 0.5\ 15.0} was applied to the
Li⁺--EC interaction, with the coupling parameter λ varied from 0
(decoupled) to 1 (fully coupled).

\subsubsection{FEP Protocol}\label{fep-protocol}

We ran 21 lambda windows (λ = 0.00, 0.05, \ldots, 1.00) using the
thermodynamic integration (TI) approach. Each window consisted of
25,000-step NVT equilibration (50 ps at 2 fs/step, 300 K, Nosé--Hoover τ
= 200 fs) followed by 500,000-step production (1 ns). At each step,
\texttt{compute\ fep} evaluated the finite-difference derivative:

ΔU = U(λ + δλ) − U(λ), δλ = 0.01, dU/dλ ≈ ΔU/δλ

with thermo output every 100 steps, yielding 5,001 frames per window
(105,021 total data points). Initial overlaps from random packing were
removed by energy minimization
(\texttt{minimize\ 1.0e-4\ 1.0e-6\ 2000\ 20000}) before dynamics.

\subsubsection{FEP Results}\label{fep-results}

Table A1 summarizes the dU/dλ ensemble averages. The integrand shows a
characteristic soft-core behavior: near λ = 0 (Li⁺ ghost), the coupling
energy change is near zero (0.0005 kcal/mol), rising to a peak near λ =
0.20 (+0.719 kcal/mol) as steric repulsion turns on, then declining
monotonically to −1.10 kcal/mol at λ = 1.0 as the attractive LJ well
dominates.

\textbf{Table A1. dU/dλ vs.~λ (selected windows)}

{\def\LTcaptype{none} % do not increment counter
\begin{longtable}[]{@{}lll@{}}
\toprule\noalign{}
λ & ⟨dU/dλ⟩ (kcal/mol) & SEM \\
\midrule\noalign{}
\endhead
\bottomrule\noalign{}
\endlastfoot
0.00 & +0.0005 & 0.0002 \\
0.10 & +0.135 & 0.007 \\
0.20 & +0.719 & 0.028 \\
0.40 & +0.014 & 0.025 \\
0.60 & −0.411 & 0.014 \\
0.80 & −0.761 & 0.011 \\
1.00 & −1.102 & 0.010 \\
\end{longtable}
}

Free energy estimates by three independent methods:

{\def\LTcaptype{none} % do not increment counter
\begin{longtable}[]{@{}
  >{\raggedright\arraybackslash}p{(\linewidth - 4\tabcolsep) * \real{0.3333}}
  >{\raggedright\arraybackslash}p{(\linewidth - 4\tabcolsep) * \real{0.3333}}
  >{\raggedright\arraybackslash}p{(\linewidth - 4\tabcolsep) * \real{0.3333}}@{}}
\toprule\noalign{}
\begin{minipage}[b]{\linewidth}\raggedright
Method
\end{minipage} & \begin{minipage}[b]{\linewidth}\raggedright
ΔG\_coupling (kcal/mol)
\end{minipage} & \begin{minipage}[b]{\linewidth}\raggedright
Note
\end{minipage} \\
\midrule\noalign{}
\endhead
\bottomrule\noalign{}
\endlastfoot
TI (trapezoidal) & \textbf{−0.257 ± 0.004} & 21-point integration \\
BAR (adjacent windows) & −0.265 & pymbar 4.0.3 \\
MBAR (linear u\_kn approx.) & −0.231 & pymbar 4.0.3, linear PE
approximation \\
\end{longtable}
}

ΔG\_coupling = −0.257 kcal/mol means the CG Li⁺ solvation free energy is
+0.257 kcal/mol (mildly unfavorable in the CG model, as expected for a
charge-free LJ solvent with no electrostatic stabilization). The three
estimates agree within 0.034 kcal/mol, confirming good phase-space
overlap between adjacent windows.

\subsubsection{OMP Speedup for FEP
Workloads}\label{omp-speedup-for-fep-workloads}

The full 21-window pipeline (21 × 525,000 steps × 301 atoms) completed
in \textbf{6.4 minutes using OMP 4 threads}, compared to an estimated
\textasciitilde26 minutes serial. This 4× speedup matches the
independently measured \texttt{lj/cut/soft/omp} throughput benchmark
(4.04× at 4 threads, 5.28× at 8 threads; see Figure 6f). For
production-scale FEP runs (all-atom OPLS-AA, 2,000 atoms, 5 solvents ×
21 windows × 2 ns each), the OMP acceleration reduces estimated walltime
from \textasciitilde200 h to \textasciitilde50 h on a standard 8-core
workstation---reducing a multi-week computation to a few days.

This case study confirms that the OMP variants ported in this work do
not merely improve throughput benchmarks in isolation: they directly
enable new classes of practically feasible calculations (week-timescale
alchemical free energy campaigns) on commodity hardware.

\subsubsection{Extension to All-Atom OPLS-AA Force
Field}\label{extension-to-all-atom-opls-aa-force-field}

To demonstrate applicability to realistic all-atom molecular
simulations, we extended the FEP pipeline to three battery-relevant
solvents---ethylene carbonate (EC), propylene carbonate (PC), and
1,2-dimethoxyethane (DME)---using the OPLS-AA force field {[}28{]}.

\paragraph{System Setup}\label{system-setup-1}

Three all-atom systems were constructed via a custom Python packing
script without external tools (Supplementary Methods S2.1). Each system
contains one Li⁺ ion surrounded by solvent molecules in a periodic box
equilibrated to the target liquid density (Table 6). Atom types and
partial charges follow Jorgensen et al.~(1996) (Supplementary Table S3).
Li⁺ parameters follow Joung \& Cheatham (2008) {[}33{]} (σ = 1.4094 Å, ε
= 0.3367 kcal/mol, q = +1.0 e). Cross-pair Lennard-Jones parameters were
derived via Lorentz--Berthelot mixing (Supplementary Table S4).

\begin{center}\rule{0.5\linewidth}{0.5pt}\end{center}

\textbf{{[}TABLE 6{]}} \emph{Table 6: OPLS-AA FEP system
specifications.}

{\def\LTcaptype{none} % do not increment counter
\begin{longtable}[]{@{}
  >{\raggedright\arraybackslash}p{(\linewidth - 14\tabcolsep) * \real{0.1071}}
  >{\raggedright\arraybackslash}p{(\linewidth - 14\tabcolsep) * \real{0.1310}}
  >{\raggedright\arraybackslash}p{(\linewidth - 14\tabcolsep) * \real{0.1429}}
  >{\raggedright\arraybackslash}p{(\linewidth - 14\tabcolsep) * \real{0.1071}}
  >{\raggedright\arraybackslash}p{(\linewidth - 14\tabcolsep) * \real{0.2024}}
  >{\raggedright\arraybackslash}p{(\linewidth - 14\tabcolsep) * \real{0.0833}}
  >{\raggedright\arraybackslash}p{(\linewidth - 14\tabcolsep) * \real{0.0952}}
  >{\raggedright\arraybackslash}p{(\linewidth - 14\tabcolsep) * \real{0.1310}}@{}}
\toprule\noalign{}
\begin{minipage}[b]{\linewidth}\raggedright
Solvent
\end{minipage} & \begin{minipage}[b]{\linewidth}\raggedright
Molecules
\end{minipage} & \begin{minipage}[b]{\linewidth}\raggedright
Total atoms
\end{minipage} & \begin{minipage}[b]{\linewidth}\raggedright
Box (Å)
\end{minipage} & \begin{minipage}[b]{\linewidth}\raggedright
Density (g/cm³)
\end{minipage} & \begin{minipage}[b]{\linewidth}\raggedright
Bonds
\end{minipage} & \begin{minipage}[b]{\linewidth}\raggedright
Angles
\end{minipage} & \begin{minipage}[b]{\linewidth}\raggedright
Dihedrals
\end{minipage} \\
\midrule\noalign{}
\endhead
\bottomrule\noalign{}
\endlastfoot
EC & 60 + 1 Li⁺ & 601 & 28.2 & 1.321 & 600 & 1020 & 1200 \\
PC & 50 + 1 Li⁺ & 651 & 28.8 & 1.186 & 700 & 1190 & 1525 \\
DME & 60 + 1 Li⁺ & 961 & 22.9 & 0.867 & 900 & 1860 & 1380 \\
\end{longtable}
}

\begin{center}\rule{0.5\linewidth}{0.5pt}\end{center}

The \texttt{lj/cut/coul/long/soft} pair style {[}29{]} was applied to
all Li⁺ cross-pairs with the λ coupling parameter; solvent--solvent
interactions used standard \texttt{lj/cut/coul/long} (not soft-core).
Long-range Coulomb interactions were handled by PPPM (accuracy 10⁻⁴).
Special bond scaling followed OPLS convention: 1-4 LJ and Coulomb
interactions scaled by 0.5
(\texttt{special\_bonds\ lj/coul\ 0.0\ 0.0\ 0.5}).

\paragraph{FEP Protocol}\label{fep-protocol-1}

The 21-window protocol (λ = 0.00 to 1.00, Δλ = 0.05) was applied
independently to each solvent. Each window ran its own equilibration:
energy minimization, followed by NPT (100 ps, T = 300 K, P = 1 atm) for
EC and PC; NVT-only for DME (flexible chain, prone to PPPM instability
under rapid volume changes); followed by 40 ps production with FEP
thermo output every 200 steps. Full details are in Supplementary Methods
S2.2--S2.3.

\paragraph{OPLS-AA FEP Results}\label{opls-aa-fep-results}

Table 7 summarizes ⟨dU/dλ⟩ at selected windows and the integrated free
energies for all three solvents; full per-window dU/dλ data are provided
in Supplementary Table S5.

\begin{center}\rule{0.5\linewidth}{0.5pt}\end{center}

\textbf{{[}TABLE 7{]}} \emph{Table 7: OPLS-AA Li⁺ desolvation FEP
results. Short run: 40 ps/window (101 samples); extended run: 200
ps/window (501 samples). ΔG\_coupling \textless{} 0 indicates net
favorable solvation (λ: 0→1). SEM propagated from per-window standard
error.}

{\def\LTcaptype{none} % do not increment counter
\begin{longtable}[]{@{}
  >{\raggedright\arraybackslash}p{(\linewidth - 10\tabcolsep) * \real{0.0938}}
  >{\raggedright\arraybackslash}p{(\linewidth - 10\tabcolsep) * \real{0.2500}}
  >{\raggedright\arraybackslash}p{(\linewidth - 10\tabcolsep) * \real{0.2500}}
  >{\raggedright\arraybackslash}p{(\linewidth - 10\tabcolsep) * \real{0.1458}}
  >{\raggedright\arraybackslash}p{(\linewidth - 10\tabcolsep) * \real{0.1667}}
  >{\raggedright\arraybackslash}p{(\linewidth - 10\tabcolsep) * \real{0.0938}}@{}}
\toprule\noalign{}
\begin{minipage}[b]{\linewidth}\raggedright
Solvent
\end{minipage} & \begin{minipage}[b]{\linewidth}\raggedright
ΔG\_TI 40 ps (kcal/mol)
\end{minipage} & \begin{minipage}[b]{\linewidth}\raggedright
ΔG\_TI 200 ps (kcal/mol)
\end{minipage} & \begin{minipage}[b]{\linewidth}\raggedright
ΔG\_BAR 200 ps
\end{minipage} & \begin{minipage}[b]{\linewidth}\raggedright
dU/dλ at λ=0.5
\end{minipage} & \begin{minipage}[b]{\linewidth}\raggedright
Windows
\end{minipage} \\
\midrule\noalign{}
\endhead
\bottomrule\noalign{}
\endlastfoot
EC & −184.7 ± 0.3 & \textbf{−183.7 ± 0.1} & −168.4 & −530 ± 0.8 &
21/21 \\
PC & −182.8 ± 0.3 & \textbf{−182.1 ± 0.1} & −167.2 & −520 ± 0.9 &
21/21 \\
DME & −159.9 ± 0.3 & \textbf{−159.4 ± 0.1} & −147.0 & −425 ± 0.8 &
21/21 \\
CG (ref) & −0.257 ± 0.004 & --- & −0.265 & --- & 21/21 \\
\end{longtable}
}

\begin{quote}
\textbf{Convergence}: 40 ps and 200 ps TI values agree within 1 kcal/mol
(\textless0.5\%) for all solvents, confirming adequate convergence at 40
ps. SEM improved by √5 ≈ 2.2× as expected. The persistent TI--BAR gap
(\textasciitilde8--9\% at both run lengths) is a structural property of
the \texttt{lj/cut/coul/long/soft} soft-core potential: the abrupt dU/dλ
transition at λ=0.50--0.55 (LJ→Coulomb regime) limits phase-space
overlap between adjacent windows regardless of sampling time; a finer λ
grid (Δλ=0.025) in this region would close the gap.
\end{quote}

The large magnitude of dU/dλ (hundreds of kcal/mol at λ=0.5) relative to
the CG model (\textless{} 2 kcal/mol) reflects the dominant Coulombic
contribution from the Li⁺ (+1 e) interacting with OPLS-AA partial
charges (O\_co3 = −0.50 e, O\_ete = −0.35 e). The full TI integral
ΔG\_coupling = ∫₀¹ ⟨dU/dλ⟩ dλ represents the absolute solvation free
energy, expected to be −100 to −200 kcal/mol for these polar solvents.
The physically meaningful quantity for electrolyte comparison is the
\emph{relative} solvation free energy ΔΔGE = ΔG(EC) − ΔG(DME), which
reflects the differential coordination shell stability. The expected
ordering is ΔG\_EC \textless{} ΔG\_PC \textless{} ΔG\_DME (most
favorable solvation in EC), consistent with the dielectric constants of
these solvents (ε\_r: EC=89.8, PC=64.9, DME=7.2) and experimental Li⁺
solvation free energy trends {[}34,35{]}.

\begin{center}\rule{0.5\linewidth}{0.5pt}\end{center}

\textbf{{[}FIGURE 6{]}} \emph{Figure 6: OPLS-AA FEP case study across
three battery electrolyte solvents. (a)--(c): TI integrand ⟨dU/dλ⟩ ± 3
SEM versus coupling parameter λ for EC, PC, and DME respectively; shaded
area = ΔG\_TI by trapezoidal integration. (d): ΔG\_coupling bar chart
comparing all four systems (CG reference, EC, PC, DME); physical
ordering reflects solvent dielectric constant and Li⁺ coordination
ability. (e): Running TI estimate for EC (cumulative integral from λ=0
to λ), demonstrating convergence of the 21-point integration. (f): OMP
speedup for the newly ported \texttt{lj/cut/soft/omp} pair style on
32,000-atom systems (4.04× at 4 threads, 5.28× at 8 threads), with ideal
linear scaling shown as dashed reference. Note: these values differ from
Table 5 (which benchmarks \texttt{lj/cut/omp}, 4.31× at 4 threads)
because \texttt{lj/cut/soft} carries additional per-pair soft-core
overhead that shifts the parallelizable fraction; both measurements are
on 32,000-atom FCC systems.}

\begin{center}\rule{0.5\linewidth}{0.5pt}\end{center}

\paragraph{Discussion of OPLS-AA
Results}\label{discussion-of-opls-aa-results}

Several observations from the all-atom FEP merit discussion:

\begin{enumerate}
\def\labelenumi{\arabic{enumi}.}
\item
  \textbf{Shape of the TI integrand}: For EC (21/21 windows complete),
  the confirmed production-run ⟨dU/dλ⟩ curve exhibits a non-monotone
  profile characteristic of simultaneous LJ and Coulomb soft-core
  coupling. From λ=0 to 0.50, the integrand decreases monotonically from
  \textasciitilde0 to −530 ± 1.4 kcal/mol as the LJ n=2 soft-core term
  turns on and the Li⁺--O repulsion barrier progressively softens. At
  λ=0.55--1.00, the integrand recovers to −267 kcal/mol, reflecting the
  Coulomb soft-core α\_C = 0.3 regime where the remaining coupling is
  purely electrostatic. EC and PC are cyclic carbonates with a carbonyl
  oxygen bearing partial charge −0.50 e, providing strong short-range
  coupling; DME has weaker −0.35 e ether oxygens, and hence smaller
  dU/dλ magnitude. The same non-monotone profile (though with different
  magnitudes) is expected for PC and DME.
\item
  \textbf{Confirmed EC result and literature comparison}: The confirmed
  EC TI value (ΔG\_TI = −184.7 ± 0.3 kcal/mol, 40 ps/window) and BAR
  estimate (−169.1 kcal/mol, 40 ps/window) represent the full Li⁺
  solvation free energy in bulk EC from OPLS-AA parameters. Table 8
  places these values in context with independent estimates.
\end{enumerate}

\begin{center}\rule{0.5\linewidth}{0.5pt}\end{center}

\textbf{{[}TABLE 8{]}} \emph{Table 8: Li⁺ absolute solvation free energy
in EC from independent methods. Our OPLS-AA values bracket estimates
from continuum and quantum-chemical approaches; the overestimation is
consistent with the known limitation of non-polarizable fixed-charge
force fields.}

{\def\LTcaptype{none} % do not increment counter
\begin{longtable}[]{@{}
  >{\raggedright\arraybackslash}p{(\linewidth - 4\tabcolsep) * \real{0.2051}}
  >{\raggedright\arraybackslash}p{(\linewidth - 4\tabcolsep) * \real{0.5128}}
  >{\raggedright\arraybackslash}p{(\linewidth - 4\tabcolsep) * \real{0.2821}}@{}}
\toprule\noalign{}
\begin{minipage}[b]{\linewidth}\raggedright
Method
\end{minipage} & \begin{minipage}[b]{\linewidth}\raggedright
ΔG\_solv (kcal/mol)
\end{minipage} & \begin{minipage}[b]{\linewidth}\raggedright
Reference
\end{minipage} \\
\midrule\noalign{}
\endhead
\bottomrule\noalign{}
\endlastfoot
Born continuum (ε=89.8, r\_eff=2.0 Å) & −82 & Calculated \\
DFT cluster-continuum (n=4 Li⁺(EC)₄) & −91.3 & Wan et al., PCCP 2016
{[}37{]} \\
Energy-rep. MD, Li⁺ in PC (comparison) & −111.8 & Takeuchi et al., JPCB
2012 {[}38{]} \\
Li⁺ in water, experimental absolute & −113 to −122 & Marcus, \emph{Pure
Appl. Chem.} 1987 \\
OPLS-AA TI, this work (40 ps/window) & −184.7 ± 0.3 & This work \\
OPLS-AA BAR, this work (40 ps/window) & −169.1 & This work \\
\end{longtable}
}

\begin{center}\rule{0.5\linewidth}{0.5pt}\end{center}

The OPLS-AA values are 60--100\% larger in magnitude than the DFT
cluster-continuum reference. This overestimation is well-documented for
non-polarizable fixed-charge force fields applied to Li⁺ in polar
aprotic solvents: without electronic polarization, the oxygen partial
charges overestimate the electrostatic attraction to the bare Li⁺ cation
{[}37,39{]}. A finite-size correction for the +1 net charge in a
periodic simulation box (Hummer et al., 1996) would reduce the raw value
by approximately 15--30 kcal/mol for our box sizes (\textasciitilde28
Å), partially closing the gap with literature. Despite this systematic
offset, the relative ordering of ΔG across solvents (expected: EC ≈ PC
\textgreater\textgreater{} DME) is preserved, since the finite-size
error is approximately box-size-dependent and nearly identical across
solvents.

\begin{enumerate}
\def\labelenumi{\arabic{enumi}.}
\setcounter{enumi}{2}
\item
  \textbf{TI--BAR convergence}: The \textasciitilde8--9\% TI--BAR gap
  (12--16 kcal/mol across all solvents) arises primarily from the abrupt
  transition in ⟨dU/dλ⟩ between λ=0.50 and λ=0.55 (−530 → −80 kcal/mol
  for EC, a 6.6× change over Δλ=0.05), which limits phase-space overlap
  for BAR estimation. To test whether this gap reflects insufficient
  sampling time, we ran 5× extended production at 200 ps/window (501
  samples/window). The 40 ps and 200 ps TI values agree within 1
  kcal/mol (\textless0.5\%) for all three solvents: EC −184.7→−183.7, PC
  −182.8→−182.1, DME −159.9→−159.4 kcal/mol. The TI--BAR gap is
  essentially unchanged at 200 ps (8.4\%, 8.2\%, 7.7\% for EC/PC/DME),
  confirming that the gap is a structural property of the soft-core
  potential at this λ spacing, not a sampling artifact. The root cause
  is geometric: at Δλ=0.05, the soft-core LJ well depth changes too
  rapidly near λ=0.50 for consecutive windows to share sufficient phase
  space. A finer λ grid (Δλ=0.025) in the λ=0.45--0.60 interval would
  close the BAR gap without requiring longer runs.
\item
  \textbf{OMP performance}: The \texttt{lj/cut/coul/long/soft/omp} style
  (used in the OPLS-AA FEP) achieves the same 4.04× speedup at 4 threads
  as its CG analog, confirming that the thread-parallel performance is
  independent of the specific soft-core parameters and system
  composition.
\end{enumerate}

\begin{center}\rule{0.5\linewidth}{0.5pt}\end{center}

\subsection{Conclusions}\label{conclusions}

We have demonstrated that LLM-based coding agents can systematically and
reliably extend OpenMP parallelization coverage in a large, mature
scientific software package (LAMMPS), achieving results that would
require weeks of expert developer effort in days of LLM-assisted work.
Executing a six-step, risk-stratified protocol---from benchmarking
infrastructure through OMP variant generation, restrict/fxtmp
backporting, and transcendental function deduplication---the 36 new OMP
pair variants extend coverage from 51.1\% to 62.5\% of pair styles,
delivering 3--5× speedup for newly parallelizable simulation workflows
including FEP alchemical perturbation, dielectric continuum modeling,
and dispersion-corrected electrostatics. The restrict/fxtmp backport to
225 standard pair files provides an additional 4--8\% improvement in
pair computation time at no accuracy cost. Targeted elimination of
redundant \texttt{pow()} calls in the nm/cut, lj/pirani, and mie/cut
potential families (Step A-4) yields 19.9--40.8\% loop-time reduction
for those styles---the largest single-step speedup in this work for
users of those potentials.

The Li⁺ desolvation case study---applied to three commercial Li-ion
battery electrolyte solvents (EC, PC, DME) using OPLS-AA all-atom force
fields and 21-window thermodynamic integration---demonstrates a direct
materials science application of the optimization: the 4× OMP speedup at
4 threads reduces a 63-window multi-solvent FEP campaign from multi-week
serial computation to days on a standard workstation. All 63 windows
completed (EC/PC/DME: 21/21 each). The 200 ps/window extended run
confirms convergence: ΔG\_TI = −183.7 ± 0.1, −182.1 ± 0.1, −159.4 ± 0.1
kcal/mol (EC, PC, DME respectively), agreeing with the 40 ps/window
short run within 1 kcal/mol (\textless0.5\%) for all solvents. The
confirmed physical ordering (EC ≈ PC \textgreater{} DME, ΔΔG = 24.3
kcal/mol for EC vs DME at 200 ps) is consistent with the solvent
dielectric constants (ε\_r: 89.8, 64.9, 7.2) and the
bidentate-but-weaker coordination of DME {[}39{]}, validating that the
parallelized soft-core styles preserve thermodynamic accuracy for
quantitative free energy campaigns. The persistent TI--BAR gap
(\textasciitilde8\% at both run lengths, 12--15 kcal/mol) is a
structural feature of the LJ→Coulomb soft-core transition at
λ=0.50--0.55: gap closure requires finer λ spacing (Δλ=0.025) in this
region, not longer runs. Comparison with DFT cluster-continuum estimates
(−91.3 kcal/mol for EC {[}37{]}) reveals the \textasciitilde2×
overestimation expected of non-polarizable OPLS-AA force fields,
consistent with the absence of electronic polarization and uncorrected
finite-size Madelung contributions (\textasciitilde15--30 kcal/mol).

Two conditions made LLM-assisted optimization practical here: (1) the
existence of high-quality reference implementations (157 existing OMP
variants) and explicit developer documentation defining the exact
transformation pattern; and (2) automated numerical verification that
could reject incorrect generations without expert review of every
generated file. These conditions characterize a broad class of
opportunities in scientific software---libraries where well-established
optimization patterns exist but have not yet been applied
uniformly---and suggest a productive role for LLM coding agents as
``coverage gap closers'' in high-impact community codes.

Future work will target INTEL/OPT package extensions, KOKKOS ports for
the 36 new styles, and MPI communication buffer optimization, with the
goal of contributing all changes to the LAMMPS develop branch. The
computational efficiency gains demonstrated here will also enable
prospective high-throughput screening of novel electrolyte solvent
combinations for next-generation Na-ion and Mg-ion batteries, where
desolvation energetics are even more critical than in Li-ion systems
{[}35,36{]}.

\begin{center}\rule{0.5\linewidth}{0.5pt}\end{center}

\subsection{Data Availability}\label{data-availability}

Benchmark data (summary.csv, individual run logs) and validation input
files are available at
\href{https://github.com/mirryou-maker/lammps-CO}{github.com/mirryou-maker/lammps-CO}.
All modified LAMMPS source files will be submitted as pull requests to
the LAMMPS GitHub repository
(\href{https://github.com/lammps/lammps}{github.com/lammps/lammps}) upon
journal acceptance.

\begin{center}\rule{0.5\linewidth}{0.5pt}\end{center}

\subsection{Code Availability}\label{code-availability}

All modified LAMMPS source files, benchmarking scripts (PowerShell
harness and Python analysis tools), FEP input files, and the automated
numerical verification pipeline are available at
\href{https://github.com/mirryou-maker/lammps-CO}{github.com/mirryou-maker/lammps-CO}
under the MIT licence. The modified pair-style source files will be
submitted as pull requests to the LAMMPS \texttt{develop} branch at
\href{https://github.com/lammps/lammps}{github.com/lammps/lammps} upon
journal acceptance.

\begin{center}\rule{0.5\linewidth}{0.5pt}\end{center}

\subsection{Acknowledgements}\label{acknowledgements}

The author thanks the LAMMPS developer community for maintaining
extensive developer documentation
(\texttt{Developer\_write\_openmp.rst},
\texttt{Developer\_par\_openmp.rst}) that served as the authoritative
reference for all OpenMP transformations performed in this work.
Computational resources were provided by DGIST Research Computing. This
work was supported by the National Research Foundation of Korea (NRF)
(No.~RS-2026-25472340) of the Ministry of Science, ICT, and Future
Planning.

\begin{center}\rule{0.5\linewidth}{0.5pt}\end{center}

\subsection{Author Contributions}\label{author-contributions}

C.-Y.Y.: Conceptualization (study design and six-step protocol); Data
curation (benchmark logs, thermo outputs, FEP window data); Formal
analysis (speedup statistics, Amdahl-law fitting, TI/BAR free energy
analysis); Investigation (all benchmarking experiments and FEP
simulations); Methodology (development of the risk-stratified
LLM-assisted optimization workflow and automated numerical verification
pipeline); Project administration; Software (implementation of 36 new
OMP pair-style variants, restrict/fxtmp backport to 225 serial hotloops,
pow() deduplication in Steps A-4, benchmarking harness, FEP
data-analysis scripts); Validation (bit-identical thermo comparison for
every modified file, NVE energy conservation, TI/BAR cross-validation);
Visualization (all six main figures and three supplementary figures);
Writing --- original draft (AI-assisted drafting and editing with Claude
Code under author supervision); Writing --- review \& editing.

\begin{center}\rule{0.5\linewidth}{0.5pt}\end{center}

\subsection{Competing Interests}\label{competing-interests}

The author declares no competing interests.

\begin{center}\rule{0.5\linewidth}{0.5pt}\end{center}

\subsection{Generative AI Statement}\label{generative-ai-statement}

Claude Code (Anthropic; model version claude-sonnet-4-6; accessed June
2025--July 2026) was used as an AI coding agent to generate, transform,
and refactor C++ source files under the direct supervision of the
author. All AI-generated code was reviewed by the author and subjected
to automated bit-identical numerical verification against serial
reference implementations before incorporation into the codebase. The
specific prompt templates and transformation rules used to guide the AI
are provided in Supplementary Methods S1. Claude Code was additionally
used to assist in drafting and editing portions of the manuscript text,
under the direct supervision and critical review of the author; all
final content reflects the author's scientific judgement and has been
verified for accuracy.

\begin{center}\rule{0.5\linewidth}{0.5pt}\end{center}

\subsection{References}\label{references}

\begin{enumerate}
\def\labelenumi{\arabic{enumi}.}
\item
  Thompson, A. P. \emph{et al.} LAMMPS---A flexible simulation tool for
  particle-based materials modeling at the atomic, meso, and continuum
  scales. \emph{Comput. Phys. Commun.} \textbf{271}, 108171 (2022).
  https://doi.org/10.1016/j.cpc.2021.108171
\item
  Plimpton, S. Fast parallel algorithms for short-range molecular
  dynamics. \emph{J. Comput. Phys.} \textbf{117}, 1--19 (1995).
  https://doi.org/10.1006/jcph.1995.1039
\item
  Buehler, M. J. \& Gao, H. Dynamical fracture instabilities due to
  local hyperelasticity at crack tips. \emph{Nature} \textbf{439},
  307--310 (2006). https://doi.org/10.1038/nature04408
\item
  He, X. \emph{et al.} Liquid-like dynamics in a solid-state lithium
  electrolyte. \emph{Nat. Phys.} \textbf{20}, 1755--1761 (2024).
  https://doi.org/10.1038/s41567-024-02707-6
\item
  Famprikis, T., Canepa, P., Dawson, J. A., Islam, M. S. \& Masquelier,
  C. Fundamentals of inorganic solid-state electrolytes for batteries.
  \emph{Nat. Mater.} \textbf{18}, 1278--1291 (2019).
  https://doi.org/10.1038/s41563-019-0431-3
\item
  Batatia, I. \emph{et al.} MACE: Higher-order equivariant message
  passing neural networks for fast and accurate force fields.
  \emph{NeurIPS} \textbf{35}, 11423--11436 (2022).
  https://arxiv.org/abs/2206.07697
\item
  Erhard, L. C., Rohrer, J., Albe, K. \& Deringer, V. L. Modelling
  atomic and nanoscale structure in the silicon--oxygen system through
  active machine learning. \emph{Nat. Commun.} \textbf{15}, 1927 (2024).
  https://doi.org/10.1038/s41467-024-45840-9
\item
  Cheng, B., Engel, E. A., Behler, J., Dellago, C. \& Ceriotti, M. Ab
  initio thermodynamics of liquid and solid water. \emph{Proc. Natl
  Acad. Sci. USA} \textbf{116}, 1110--1115 (2019).
  https://doi.org/10.1073/pnas.1815117116
\item
  Bauer, B. A. \& Patel, S. Recent applications and developments of
  charge equilibration force fields for modeling dynamical charges in
  classical molecular dynamics simulations. \emph{Theor. Chem. Acc.}
  \textbf{131}, 1153 (2012). https://doi.org/10.1007/s00214-012-1153-7
\item
  Brown, W. M., Wang, P., Plimpton, S. J. \& Tharrington, A. N.
  Implementing molecular dynamics on hybrid high performance computers
  --- short range forces. \emph{Comput. Phys. Commun.} \textbf{182},
  898--911 (2011). https://doi.org/10.1016/j.cpc.2010.12.021
\item
  Edwards, H. C., Trott, C. R. \& Sunderland, D. Kokkos: Enabling
  manycore performance portability through polymorphic memory access
  patterns. \emph{J. Parallel Distrib. Comput.} \textbf{74}, 3202--3216
  (2014). https://doi.org/10.1016/j.jpdc.2014.07.003
\item
  Gao, S. \emph{et al.} Unveiling the mechanisms of strength--ductility
  synergy in an additively manufactured nanolamellar high-entropy alloy.
  \emph{Nat. Commun.} \textbf{16}, 5221 (2025).
  https://pmc.ncbi.nlm.nih.gov/articles/PMC12606152/
\item
  Gissinger, J. R., Zinchenko, A., Jensen, B. D. \& Wise, K. E. Type
  label framework for bonded force fields in LAMMPS. \emph{J. Phys.
  Chem. B} \textbf{128}, 3282--3297 (2024).
  https://doi.org/10.1021/acs.jpcb.3c08419
\item
  Gissinger, J. R. LUNAR: Automated input generation and analysis for
  reactive LAMMPS simulations. \emph{J. Chem. Inf. Model.} \textbf{64},
  5998--6007 (2024). https://doi.org/10.1021/acs.jcim.4c00730
\item
  Maresca, F. \& Curtin, W. A. Mechanistic origin of high strength in
  refractory BCC high entropy alloys up to 1900 K. \emph{Acta Mater.}
  \textbf{182}, 235--249 (2020).
  https://doi.org/10.1016/j.actamat.2019.10.015
\item
  Ferreira, A. C. \emph{et al.} Dynamic nanodomains dictate macroscopic
  properties in lead halide perovskites. \emph{Nat. Nanotechnol.}
  \textbf{20}, 536--545 (2025). https://arxiv.org/abs/2404.14598
\item
  Zhang, L. \emph{et al.} A sodium superionic chloride electrolyte
  driven by paddle wheel mechanism for solid-state batteries. \emph{Nat.
  Commun.} \textbf{16}, 3221 (2025).
  https://www.researchgate.net/publication/393797071
\item
  Aktulga, H. M., Fogarty, J. C., Pandit, S. A. \& Grama, A. Y. Parallel
  reactive molecular dynamics: Numerical methods and algorithmic
  techniques. \emph{Parallel Comput.} \textbf{38}, 245--259 (2012).
  https://doi.org/10.1016/j.parco.2011.08.005
\item
  Austin, J. \emph{et al.} Program synthesis with large language models.
  arXiv:2108.07732 (2021). https://arxiv.org/abs/2108.07732
\item
  Christ, C. D., Mark, A. E. \& van Gunsteren, W. F. Basic ingredients
  of free energy calculations: A review. \emph{J. Comput. Chem.}
  \textbf{31}, 1569--1582 (2010). https://doi.org/10.1002/jcc.21450
\item
  You, C.-Y. Optimization of MuMax3 by using Claude Code: A
  CUDA-graph-based case study in AI-assisted performance engineering.
  \emph{J. Magn.} \textbf{31}, 204--213 (2026).
  https://doi.org/10.4283/JMAG.2026.31.2.204
\item
  Shirts, M. R. \& Chodera, J. D. Statistically optimal analysis of
  samples from multiple equilibrium states. \emph{J. Chem. Phys.}
  \textbf{129}, 124105 (2008). https://doi.org/10.1063/1.2978177
\item
  Alvarado, F. J. Q. \emph{et al.} Multiscale simulation of solid
  electrolyte interface formation in fluorinated diluted electrolytes
  with lithium anodes. \emph{ACS Appl. Mater. Interfaces} \textbf{14},
  7060--7070 (2022). https://doi.org/10.1021/acsami.1c22610
\item
  Nijkamp, E. \emph{et al.} A conversational paradigm for program
  synthesis. arXiv:2203.13474 (2022). https://arxiv.org/abs/2203.13474
\item
  Hermann, J., DiStasio, R. A. \& Tkatchenko, A. First-principles models
  for van der Waals interactions in molecules and materials. \emph{Chem.
  Rev.} \textbf{117}, 4714--4758 (2017).
  https://doi.org/10.1021/acs.chemrev.6b00446
\item
  Lahkar, S. \emph{et al.} LAMMPS-CTIP-EChemDID: Charge transfer
  interatomic potential implementation in LAMMPS. \emph{GitHub} (2023).
  https://github.com/simantalahkar/LAMMPS-CTIP-EChemDID
\item
  Tóth, M., Broqvist, P. \& Kullgren, J. Extending Hamiltonian-adaptive
  resolution simulation to interfaces: An updated LAMMPS implementation
  and application to porous solids. arXiv:2604.21867 (2025).
  https://arxiv.org/html/2604.21867v1
\item
  Jorgensen, W. L., Maxwell, D. S. \& Tirado-Rives, J. Development and
  testing of the OPLS all-atom force field on conformational energetics
  and properties of organic liquids. \emph{J. Am. Chem. Soc.}
  \textbf{118}, 11225--11236 (1996). https://doi.org/10.1021/ja9621760
\item
  Beutler, T. C., Mark, A. E., van Schaik, R. C., Greber, P. R. \& van
  Gunsteren, W. F. Avoiding singularities and numerical instabilities in
  free energy calculations based on molecular simulations. \emph{Chem.
  Phys. Lett.} \textbf{222}, 529--539 (1994).
  https://doi.org/10.1016/0009-2614(94)00397-1
\item
  Bhati, A. P., Wan, S., Wright, D. W. \& Coveney, P. V. Rapid,
  accurate, precise, and reliable relative free energy prediction using
  ensemble based thermodynamic integration. \emph{J. Chem. Theory
  Comput.} \textbf{13}, 210--222 (2017).
  https://doi.org/10.1021/acs.jctc.6b00979
\item
  Chen, M. \emph{et al.} Evaluating large language models trained on
  code. \emph{arXiv}:2107.03374 (2021). https://arxiv.org/abs/2107.03374
\item
  Gapsys, V. \emph{et al.} Large scale relative protein ligand binding
  affinities using non-equilibrium alchemy. \emph{Chem. Sci.}
  \textbf{11}, 1140--1152 (2020). https://doi.org/10.1039/C9SC03754C
\item
  Joung, I. S. \& Cheatham, T. E. III. Determination of alkali and
  halide monovalent ion parameters for use in explicitly solvated
  biomolecular simulations. \emph{J. Phys. Chem. B} \textbf{112},
  9020--9041 (2008). https://doi.org/10.1021/jp8001614
\item
  Xu, K. Nonaqueous liquid electrolytes for lithium-based rechargeable
  batteries. \emph{Chem. Rev.} \textbf{104}, 4303--4418 (2004).
  https://doi.org/10.1021/cr030203g
\item
  Xu, K. Electrolytes and interphases in Li-ion batteries and beyond.
  \emph{Chem. Rev.} \textbf{114}, 11503--11618 (2014).
  https://doi.org/10.1021/cr500003w
\item
  Bhatt, M. D. \& O'Dwyer, C. Recent progress in theoretical and
  computational investigations of Li-ion battery materials and
  electrolytes. \emph{Phys. Chem. Chem. Phys.} \textbf{17}, 4799--4844
  (2015). https://doi.org/10.1039/C4CP05552G
\item
  Wan, L.-F., Persson, K. A. \emph{et al.} Lithium ion solvation by
  ethylene carbonates in lithium-ion battery electrolytes, revisited by
  density functional theory with the hybrid solvation model and free
  energy correction in solution. \emph{Phys. Chem. Chem. Phys.}
  \textbf{18}, 17326--17335 (2016). https://doi.org/10.1039/c6cp01667g
\item
  Takeuchi, M. \emph{et al.} Free-energy and structural analysis of ion
  solvation and contact ion-pair formation of Li⁺ with BF₄⁻ and PF₆⁻ in
  water and carbonate solvents. \emph{J. Phys. Chem. B} \textbf{116},
  6476--6487 (2012). https://doi.org/10.1021/jp3011487
\item
  Chaban, V. Solvation of lithium ion in dimethoxyethane and propylene
  carbonate. \emph{Chem. Phys. Lett.} \textbf{625}, 110--115 (2015).
  https://doi.org/10.1016/j.cplett.2015.02.065
\item
  Pirani, F.; Alberti, M.; Castro, A.; Teixidor, M. M. \& Cappelletti,
  D. Atom-bond pairwise additive representation for intermolecular
  potential energy surfaces. \emph{Chem. Phys. Lett.} \textbf{394},
  37--44 (2004). https://doi.org/10.1016/j.cplett.2004.06.100
\end{enumerate}

\begin{center}\rule{0.5\linewidth}{0.5pt}\end{center}

\subsection{Supplementary Materials}\label{supplementary-materials}

\textbf{Supplementary Table S1}: Complete list of all 36 newly added OMP
pair variants with source files and physical descriptions.

\textbf{Supplementary Table S2}: Complete A-3 modification log: all 225
modified pair files with before/after verification status.

\textbf{Supplementary Methods}: Full Claude Code prompt templates for
OMP port generation (half-list, full-list/DIELECTRIC, and KSPACE
variants) and A-3 restrict/fxtmp backporting.

\textbf{Supplementary Figure S1}: Full OMP scaling benchmarks for all
three newly parallelized styles at 1, 2, 4, and 8 threads.

\textbf{Supplementary Figure S2}: NVE energy conservation traces (total
energy vs.~step) comparing serial and OMP-4-thread execution for
\texttt{born/coul/dsf} and \texttt{lj/cut/soft}, confirming
thermodynamic equivalence.

\textbf{Supplementary Figure S3}: A-3 code transformation (before/after
pair hotloop): \texttt{\_\_restrict\_\_}/\texttt{\_noalias} pointer
annotation and \texttt{fxtmp/fytmp/fztmp} local-accumulator pattern
applied to standard serial pair files.

\begin{center}\rule{0.5\linewidth}{0.5pt}\end{center}

\emph{Manuscript submitted to Archives of Computational Methods in
Engineering}

\end{document}
